\chapter{Heap binomiali}\label{ch:heap-binomiali}
Gli \textbf{heap binomiali} sono una struttura dati per gestire un insieme di elementi con priorità, come gli heap binari, ma permettono anche l'\textbf{unione} di due heap in tempo logaritmico. Sono stati introdotti da Jean Vuillemin nel 1978.

Un heap binomiale è una collezione di alberi con una forma molto rigida, detti alberi binomiali. Due collezioni di questo tipo si fondono con un procedimento simile alla somma di due numeri binari.

In questo capitolo partiamo dalle operazioni che vogliamo supportare, definiamo gli alberi binomiali e le loro proprietà, e da questi arriviamo agli heap binomiali e alle loro operazioni. Chiude il capitolo l'analisi ammortizzata di una sequenza di inserimenti.

\section{Mergeable heap}\label{sec:hb-mergeable}
Un \textbf{mergeable heap} (heap unificabile) è una struttura dati che contiene un insieme di nodi, ciascuno con una chiave, e che supporta le seguenti operazioni:
\begin{itemize}
    \item \MakeHeap$()$: crea e restituisce un nuovo heap vuoto.
    \item \Insert$(H,x)$: inserisce nell'heap $H$ il nodo $x$, la cui chiave è già stata assegnata.
    \item \Minimum$(H)$: restituisce un puntatore al nodo di $H$ con la chiave minima.
    \item \ExtractMin$(H)$: elimina da $H$ il nodo con la chiave minima, e ne restituisce il puntatore.
    \item \Union$(H_1,H_2)$: crea e restituisce un nuovo heap con tutti i nodi di $H_1$ e di $H_2$. I due heap di partenza vengono ``distrutti'' dall'operazione.
    \item \DecreaseKey$(H,x,k)$: assegna al nodo $x$ di $H$ la nuova chiave $k$, che si suppone non maggiore della chiave attuale.
    \item \Delete$(H,x)$: elimina il nodo $x$ da $H$.
\end{itemize}
Di queste operazioni considereremo due implementazioni:
\begin{itemize}
    \item \textbf{Heap binomiali} (J.~Vuillemin, 1978): tutte le operazioni costano $O(\lg n)$ nel caso peggiore, unione compresa.
    \item \textbf{Heap di Fibonacci} (M.\,L.~Fredman e R.\,E.~Tarjan, 1984--1987): \Insert{}, \Minimum{}, \Union{} e \DecreaseKey{} costano $\Theta(1)$ in senso ammortizzato.
\end{itemize}
La tabella~\ref{tab:hb-confronto} ne confronta i costi con quelli degli heap binari.

\begin{table}[htbp]
\centering
\small
\begin{tabular}{lccc}
\toprule
Operazione & Heap binario & Heap binomiale & Heap di Fibonacci \\
 & (caso peggiore) & (caso peggiore) & (ammortizzato) \\
\midrule
\MakeHeap{} & $\Theta(1)$ & $\Theta(1)$ & $\Theta(1)$ \\
\Insert{} & $\Theta(\lg n)$ & $O(\lg n)$ & $\Theta(1)$ \\
\Minimum{} & $\Theta(1)$ & $O(\lg n)$ & $\Theta(1)$ \\
\ExtractMin{} & $\Theta(\lg n)$ & $O(\lg n)$ & $O(\lg n)$ \\
\Union{} & $\Theta(n)$ & $O(\lg n)$ & $\Theta(1)$ \\
\DecreaseKey{} & $\Theta(\lg n)$ & $O(\lg n)$ & $\Theta(1)$ \\
\Delete{} & $\Theta(\lg n)$ & $O(\lg n)$ & $O(\lg n)$ \\
\bottomrule
\end{tabular}
\caption{Il costo delle operazioni in tre implementazioni dei mergeable heap, dove $n$ è il numero di elementi. Per gli heap binari e binomiali il costo è quello del caso peggiore, mentre per gli heap di Fibonacci è un costo ammortizzato.}
\label{tab:hb-confronto}
\end{table}

\section{Alberi binomiali}\label{sec:hb-alberi}
Un heap binomiale è una collezione di alberi con una forma molto particolare, detti alberi binomiali. Prima di vedere gli heap, quindi, studiamo questi alberi e le loro proprietà.

\begin{definizione}[Albero binomiale]\label{def:albero-binomiale}
Per ogni $k \in \mathbb{N}$, l'\textbf{albero binomiale} $B_k$ di grado $k$ è l'albero ordinato definito dalla seguente ricorsione:
\begin{itemize}
    \item $B_0$ è l'albero formato da un solo nodo.
    \item Per $k \ge 1$, $B_k$ si ottiene da due copie di $B_{k-1}$: la radice di una delle due diventa il figlio più a sinistra della radice dell'altra.
\end{itemize}
\end{definizione}

Gli alberi binomiali sono \textbf{alberi ordinati}, ovvero alberi in cui i figli di ogni nodo hanno un ordine: si parla di primo figlio, secondo figlio, e così via. Il \textbf{grado} di un nodo è il numero dei suoi figli, e il grado di un albero binomiale è quello della sua radice.

La definizione è ricorsiva: per costruire $B_k$ si prendono due copie di $B_{k-1}$, e si appende la radice di una delle due sotto la radice dell'altra, come primo figlio (figura~\ref{fig:hb-definizione}a). Ad esempio, $B_1$ è formato da due nodi collegati da un arco, $B_2$ da due copie di $B_1$, e così via.

\begin{figure}[htbp]
\centering
\begin{tikzpicture}[binheap]
  % \sotto{x}{y}{mezza base}{altezza}{etichetta}{colore}: un sottoalbero con vertice in (x,y),
  % con il nome sotto la base
  \newcommand{\sotto}[6]{%
    \draw[#6, fill=#6!8] (#1,#2) -- ({#1-#3},{#2-#4}) -- ({#1+#3},{#2-#4}) -- cycle;
    \node[font=\footnotesize] at (#1,{#2-#4-0.22}) {$#5$};}
  % (a) la definizione ricorsiva
  \node[lab, anchor=west] at (-0.7,0.45) {(a)};
  \node[vn] at (0,0) {};
  \node[lab] at (0,-2.45) {$B_0$};
  \sotto{1.45}{-0.55}{0.45}{1.0}{B_{k-1}}{FigColor}
  \sotto{2.6}{0}{0.45}{1.0}{B_{k-1}}{Accent}
  \node[vn, draw=FigColor, fill=FigColor!30] (r1) at (1.45,-0.55) {};
  \node[vn] (r2) at (2.6,0) {};
  \draw (r2) -- (r1);
  \draw[decorate, decoration={brace, mirror, amplitude=4pt}, TextGray] (0.95,-2.05) -- node[lab, below=4pt] {$B_k$} (3.1,-2.05);
  % (b) i figli della radice
  \node[lab, anchor=west] at (4.2,0.45) {(b)};
  \node[vn] (r) at (9.3,0.25) {};
  \sotto{5.6}{-0.35}{0.55}{1.25}{B_{k-1}}{Accent}
  \sotto{6.95}{-0.35}{0.45}{1.0}{B_{k-2}}{Accent}
  \node at (7.85,-0.75) {$\cdots$};
  \sotto{8.55}{-0.35}{0.28}{0.65}{B_2}{Accent}
  \sotto{9.1}{-0.35}{0.18}{0.42}{B_1}{Accent}
  \foreach \x/\n in {5.6/c1, 6.95/c2, 8.55/c3, 9.1/c4, 9.55/c5} \node[vn] (\n) at (\x,-0.35) {};
  \foreach \n in {c1, c2, c3, c4, c5} \draw (r) -- (\n);
  \node[font=\footnotesize] at (9.6,-0.67) {$B_0$};
  \draw[decorate, decoration={brace, mirror, amplitude=4pt}, TextGray] (5.0,-2.05) -- node[lab, below=4pt] {$B_k$} (9.8,-2.05);
\end{tikzpicture}
\caption{(a) La definizione ricorsiva degli alberi binomiali: $B_0$ è un solo nodo, mentre $B_k$ è formato da due copie di $B_{k-1}$, in cui la radice di una copia diventa il primo figlio della radice dell'altra. I triangoli rappresentano sottoalberi. (b) Un altro modo di vedere $B_k$: i figli della radice, da sinistra a destra, sono le radici di $B_{k-1}, B_{k-2}, \ldots, B_1, B_0$.}
\label{fig:hb-definizione}
\end{figure}

La figura~\ref{fig:hb-alberi} mostra i primi cinque alberi binomiali. Si notano subito alcune regolarità: il numero di nodi raddoppia a ogni passo, l'altezza cresce di uno, e la radice di $B_k$ ha $k$ figli. Inoltre i nodi di $B_4$ alle profondità $0, 1, 2, 3, 4$ sono $1, 4, 6, 4, 1$, ovvero una riga del triangolo di Tartaglia: è da qui che gli alberi binomiali prendono il nome.

\begin{figure}[htbp]
\centering
\begin{widget}{binomial-tree}[k=4]
\begin{tikzpicture}[binheap]
  \node[vn] (b0-r) at (0.000,0.000) {};
  \node[lab] at (0.000,-2.550) {$B_0$};
  \node[vn] (b1-0) at (0.910,-0.500) {};
  \node[vn] (b1-1) at (0.910,0.000) {};
  \draw (b1-1) -- (b1-0);
  \node[lab] at (0.910,-2.550) {$B_1$};
  \node[vn] (b2-00) at (1.820,-1.000) {};
  \node[vn] (b2-01) at (1.820,-0.500) {};
  \node[vn] (b2-10) at (2.180,-0.500) {};
  \node[vn] (b2-11) at (2.180,0.000) {};
  \draw (b2-01) -- (b2-00);
  \draw (b2-11) -- (b2-01);
  \draw (b2-11) -- (b2-10);
  \node[lab] at (2.000,-2.550) {$B_2$};
  \node[vn] (b3-000) at (3.290,-1.500) {};
  \node[vn] (b3-001) at (3.290,-1.000) {};
  \node[vn] (b3-010) at (3.650,-1.000) {};
  \node[vn] (b3-011) at (3.650,-0.500) {};
  \node[vn] (b3-100) at (4.010,-1.000) {};
  \node[vn] (b3-101) at (4.010,-0.500) {};
  \node[vn] (b3-110) at (4.370,-0.500) {};
  \node[vn] (b3-111) at (4.370,0.000) {};
  \draw (b3-001) -- (b3-000);
  \draw (b3-011) -- (b3-001);
  \draw (b3-011) -- (b3-010);
  \draw (b3-111) -- (b3-011);
  \draw (b3-101) -- (b3-100);
  \draw (b3-111) -- (b3-101);
  \draw (b3-111) -- (b3-110);
  \node[lab] at (3.830,-2.550) {$B_3$};
  \node[vn] (b4-0000) at (5.480,-2.000) {};
  \node[vn] (b4-0001) at (5.480,-1.500) {};
  \node[vn] (b4-0010) at (5.840,-1.500) {};
  \node[vn] (b4-0011) at (5.840,-1.000) {};
  \node[vn] (b4-0100) at (6.200,-1.500) {};
  \node[vn] (b4-0101) at (6.200,-1.000) {};
  \node[vn] (b4-0110) at (6.560,-1.000) {};
  \node[vn] (b4-0111) at (6.560,-0.500) {};
  \node[vn] (b4-1000) at (6.920,-1.500) {};
  \node[vn] (b4-1001) at (6.920,-1.000) {};
  \node[vn] (b4-1010) at (7.280,-1.000) {};
  \node[vn] (b4-1011) at (7.280,-0.500) {};
  \node[vn] (b4-1100) at (7.640,-1.000) {};
  \node[vn] (b4-1101) at (7.640,-0.500) {};
  \node[vn] (b4-1110) at (8.000,-0.500) {};
  \node[vn] (b4-1111) at (8.000,0.000) {};
  \draw (b4-0001) -- (b4-0000);
  \draw (b4-0011) -- (b4-0001);
  \draw (b4-0011) -- (b4-0010);
  \draw (b4-0111) -- (b4-0011);
  \draw (b4-0101) -- (b4-0100);
  \draw (b4-0111) -- (b4-0101);
  \draw (b4-0111) -- (b4-0110);
  \draw (b4-1111) -- (b4-0111);
  \draw (b4-1001) -- (b4-1000);
  \draw (b4-1011) -- (b4-1001);
  \draw (b4-1011) -- (b4-1010);
  \draw (b4-1111) -- (b4-1011);
  \draw (b4-1101) -- (b4-1100);
  \draw (b4-1111) -- (b4-1101);
  \draw (b4-1111) -- (b4-1110);
  \node[lab] at (6.740,-2.550) {$B_4$};
  \node[lab, anchor=south] at (9.010,0.22) {$i$};
  \node[lab, anchor=south west] at (9.560,0.22) {nodi};
  \draw[RuleGray, densely dotted] (8.460,0.000) -- (8.860,0.000);
  \node[lab, anchor=west] at (8.910,0.000) {$0$};
  \node[lab, anchor=west] at (9.610,0.000) {$\binom{4}{0} = 1$};
  \draw[RuleGray, densely dotted] (8.460,-0.500) -- (8.860,-0.500);
  \node[lab, anchor=west] at (8.910,-0.500) {$1$};
  \node[lab, anchor=west] at (9.610,-0.500) {$\binom{4}{1} = 4$};
  \draw[RuleGray, densely dotted] (8.460,-1.000) -- (8.860,-1.000);
  \node[lab, anchor=west] at (8.910,-1.000) {$2$};
  \node[lab, anchor=west] at (9.610,-1.000) {$\binom{4}{2} = 6$};
  \draw[RuleGray, densely dotted] (8.460,-1.500) -- (8.860,-1.500);
  \node[lab, anchor=west] at (8.910,-1.500) {$3$};
  \node[lab, anchor=west] at (9.610,-1.500) {$\binom{4}{3} = 4$};
  \draw[RuleGray, densely dotted] (8.460,-2.000) -- (8.860,-2.000);
  \node[lab, anchor=west] at (8.910,-2.000) {$4$};
  \node[lab, anchor=west] at (9.610,-2.000) {$\binom{4}{4} = 1$};
\end{tikzpicture}
\end{widget}
\caption{Gli alberi binomiali da $B_0$ a $B_4$. Ogni albero è formato da due copie del precedente, quindi il numero di nodi raddoppia e l'altezza cresce di uno. A destra il numero di nodi di $B_4$ a ciascuna profondità, dato dai coefficienti binomiali $\binom{4}{i}$.}
\label{fig:hb-alberi}
\end{figure}

\subsection{Proprietà degli alberi binomiali}\label{sec:hb-proprieta}
Il lemma~\ref{lem:hb-proprieta} raccoglie le regolarità che abbiamo osservato.

\begin{lemma}[Proprietà degli alberi binomiali]\label{lem:hb-proprieta}
Per ogni $k = 0, 1, 2, \ldots$ valgono le seguenti proprietà:
\begin{enumerate}
    \item $B_k$ ha $2^k$ nodi.
    \item L'altezza di $B_k$ è $k$.
    \item $B_k$ ha esattamente $\binom{k}{i}$ nodi a profondità $i$, per $i = 0, 1, \ldots, k$.
    \item Per i gradi dei nodi e i figli della radice si ha:
    \begin{enumerate}[label=(\alph*)]
        \item la radice di $B_k$ ha grado $k$;
        \item ogni altro nodo di $B_k$ ha grado minore di $k$;
        \item i figli della radice di $B_k$ sono, nell'ordine, le radici di $B_{k-1}, \ldots, B_1, B_0$.
    \end{enumerate}
\end{enumerate}
\end{lemma}

\paragraph{Dimostrazione.}
Procediamo per induzione su $k$. Nel caso base, $k = 0$, l'albero $B_0$ è formato da un solo nodo, quindi le quattro proprietà sono immediate. $B_0$ ha $1 = 2^0$ nodi e altezza $0$, ha $1 = \binom{0}{0}$ nodi a profondità $0$, e la sua radice, l'unico nodo, ha grado $0$.

Per il passo induttivo supponiamo che il lemma valga per $k-1$, con $k \ge 1$, e consideriamo le due copie di $B_{k-1}$ che formano $B_k$. Chiamiamo \textbf{copia di destra} quella che conserva la radice, e \textbf{copia di sinistra} quella la cui radice diventa il primo figlio della radice di $B_k$ (figura~\ref{fig:hb-dimostrazione}). Dimostriamo le quattro proprietà nell'ordine del lemma:

\begin{enumerate}
    \item Il numero di nodi di $B_k$ è la somma dei nodi delle due copie. Indicando con $|B_k|$ il numero di nodi di $B_k$, per l'ipotesi induttiva si ha:
    \[
    |B_k| = |B_{k-1}| + |B_{k-1}| = 2 \cdot 2^{k-1} = 2^k.
    \]
    \item I nodi della copia di destra mantengono la loro profondità, mentre quelli della copia di sinistra scendono di un livello. Il nodo più profondo di $B_k$, quindi, è il nodo più profondo della copia di sinistra, e l'altezza vale:
    \[
    h(B_k) = h(B_{k-1}) + 1 = (k-1) + 1 = k,
    \]
    dove $h(B)$ indica l'altezza dell'albero $B$.
    \item Sia $D(k,i)$ il numero di nodi di $B_k$ a profondità $i$. Un nodo di $B_k$ a profondità $i$ è un nodo a profondità $i$ della copia di destra, oppure un nodo a profondità $i-1$ della copia di sinistra, che è scesa di un livello. Per l'ipotesi induttiva, quindi, se $1 \le i \le k-1$ si ha:
    \[
    D(k,i) = D(k-1,i) + D(k-1,i-1) = \binom{k-1}{i} + \binom{k-1}{i-1} = \binom{k}{i},
    \]
    dove l'ultima uguaglianza è la regola con cui si costruisce il triangolo di Tartaglia. Restano i due estremi. A profondità $0$ c'è soltanto la radice, e $1 = \binom{k}{0}$. A profondità $k$ ci sono soltanto i nodi a profondità $k-1$ della copia di sinistra, ovvero $\binom{k-1}{k-1} = 1 = \binom{k}{k}$.
    \item Dimostriamo separatamente i tre punti.
    \begin{enumerate}[label=(\alph*)]
        \item La radice di $B_k$ è la radice della copia di destra, che riceve un figlio in più: il suo grado vale $(k-1) + 1 = k$.
        \item Ogni altro nodo appartiene a una delle due copie, e il suo grado non cambia, in quanto l'unico nodo che riceve un nuovo figlio è la radice. Per l'ipotesi induttiva questo grado è al più $k-1$, ovvero minore di $k$.
        \item Il primo figlio della radice di $B_k$ è la radice della copia di sinistra, ovvero di un $B_{k-1}$. I figli successivi sono quelli della radice della copia di destra, che per l'ipotesi induttiva sono le radici di $B_{k-2}, \ldots, B_1, B_0$, in quest'ordine. \hfill $\square$
    \end{enumerate}
\end{enumerate}

\begin{figure}[htbp]
\centering
\begin{tikzpicture}[binheap]
  % le profondità
  \foreach \y/\l in {0/0, -0.6/1, -1.8/i, -3.0/k} {
    \draw[RuleGray, densely dotted] (-0.2,\y) -- (4.4,\y);
    \node[lab, anchor=east] at (-0.25,\y) {$\l$};
  }
  \node[lab, anchor=east] at (-0.25,0.55) {profondità};
  % copia di destra, che conserva la radice, e copia di sinistra, scesa di un livello
  \draw[Accent, fill=Accent!8] (3.2,0) -- (2.5,-2.4) -- (3.9,-2.4) -- cycle;
  \draw[FigColor, fill=FigColor!8] (1.6,-0.6) -- (0.9,-3.0) -- (2.3,-3.0) -- cycle;
  \node[vn] (r2) at (3.2,0) {};
  \node[vn, draw=FigColor, fill=FigColor!30] (r1) at (1.6,-0.6) {};
  \draw (r2) -- (r1);
  % i nodi a profondità i delle due copie
  \draw[FigColor, line width=2.2pt] (1.25,-1.8) -- (1.95,-1.8);
  \draw[Accent, line width=2.2pt] (2.675,-1.8) -- (3.725,-1.8);
  \node[lab, font=\footnotesize, text=FigColor, anchor=east] at (0.95,-1.25) {copia di sinistra};
  \node[lab, font=\footnotesize, text=Accent, anchor=west] at (3.7,-0.95) {copia di destra};
  % i conteggi
  \node[lab, anchor=west] at (4.6,0) {$1 = \binom{k}{0}$};
  \node[lab, anchor=west] at (4.6,-1.8) {$\textcolor{FigColor}{\binom{k-1}{i-1}} + \textcolor{Accent}{\binom{k-1}{i}} = \binom{k}{i}$};
  \node[lab, anchor=west] at (4.6,-3.0) {$\binom{k-1}{k-1} = 1 = \binom{k}{k}$};
\end{tikzpicture}
\caption{Il passo induttivo della proprietà 3. La copia di destra conserva la radice e la profondità dei suoi nodi, mentre la copia di sinistra scende di un livello. A profondità $i$ si trovano quindi i $\binom{k-1}{i}$ nodi a profondità $i$ della copia di destra e i $\binom{k-1}{i-1}$ nodi a profondità $i-1$ della copia di sinistra.}
\label{fig:hb-dimostrazione}
\end{figure}

\begin{osservazione}[Perché binomiali]
La regola $\binom{k}{i} = \binom{k-1}{i} + \binom{k-1}{i-1}$ ha una lettura combinatoria che ricalca la costruzione di $B_k$. Fissato un elemento in un insieme di $k$ elementi, i sottoinsiemi di $i$ elementi si dividono in quelli che non lo contengono, che sono $\binom{k-1}{i}$, e in quelli che lo contengono, che sono $\binom{k-1}{i-1}$. Allo stesso modo, i nodi di $B_k$ a profondità $i$ si dividono tra la copia di destra e la copia di sinistra.
\end{osservazione}

Dalle proprietà 4(a) e 4(b) segue un limite al grado dei nodi, che useremo per stimare il costo delle operazioni.

\begin{corollario}[Grado massimo in un albero binomiale]\label{cor:hb-grado}
Sia $B$ un albero binomiale con $n$ nodi. Allora ogni nodo di $B$ ha grado al più $\lg n$.
\end{corollario}

\paragraph{Dimostrazione.}
Per qualche $k \in \mathbb{N}$ si ha $B = B_k$, quindi per la proprietà 1 vale $n = 2^k$, ovvero $k = \lg n$. Per le proprietà 4(a) e 4(b) ogni nodo di $B_k$ ha grado al più $k$, da cui la tesi. \hfill $\square$

\section{Heap binomiali}\label{sec:hb-heap}
Possiamo ora definire gli heap binomiali, come collezioni di alberi binomiali con due vincoli.

\begin{definizione}[Heap binomiale]\label{def:heap-binomiale}
Un \textbf{heap binomiale} $H$ è una foresta di alberi binomiali tale che:
\begin{itemize}
    \item Ogni albero binomiale di $H$ gode della proprietà di \textbf{min-heap}: la chiave di ogni nodo è maggiore o uguale alla chiave di suo padre.
    \item Per ogni $k \in \mathbb{N}$, $H$ contiene al più un albero binomiale di grado $k$.
\end{itemize}
\end{definizione}

La figura~\ref{fig:hb-esempio} mostra un heap binomiale con $13$ nodi, formato da tre alberi binomiali di grado $0$, $2$ e $3$.

\begin{figure}[htbp]
\centering
\begin{tikzpicture}[binheap, >=Stealth]
  % (a) 
  \node[hn] (s0n0) at (1.530,0.000) {10};
  \node[hn] (s0n3) at (2.510,-1.200) {18};
  \node[hn] (s0n2) at (2.510,-0.600) {12};
  \node[hn] (s0n4) at (3.070,-0.600) {25};
  \node[hn] (s0n1) at (3.070,0.000) {1};
  \node[hn] (s0n8) at (4.050,-1.800) {27};
  \node[hn] (s0n7) at (4.050,-1.200) {11};
  \node[hn] (s0n9) at (4.610,-1.200) {17};
  \node[hn] (s0n6) at (4.610,-0.600) {8};
  \node[hn] (s0n11) at (5.170,-1.200) {38};
  \node[hn] (s0n10) at (5.170,-0.600) {14};
  \node[hn] (s0n12) at (5.730,-0.600) {29};
  \node[hn] (s0n5) at (5.730,0.000) {6};
  \draw (s0n2) -- (s0n3);
  \draw (s0n1) -- (s0n2);
  \draw (s0n1) -- (s0n4);
  \draw (s0n7) -- (s0n8);
  \draw (s0n6) -- (s0n7);
  \draw (s0n6) -- (s0n9);
  \draw (s0n5) -- (s0n6);
  \draw (s0n10) -- (s0n11);
  \draw (s0n5) -- (s0n10);
  \draw (s0n5) -- (s0n12);
  \draw[->] (s0n0) -- (s0n1);
  \draw[->] (s0n1) -- (s0n5);
  \node[head, anchor=east] (s0h0) at (1.094,0.000) {$\mathit{head}[H]$};
  \draw[->] (s0h0) -- (s0n0);
\end{tikzpicture}
\caption{Un heap binomiale con $13$ nodi, formato dagli alberi binomiali $B_0$, $B_2$ e $B_3$. Le radici sono collegate nella lista delle radici, ordinata per grado crescente, e $\mathit{head}[H]$ punta alla prima.}
\label{fig:hb-esempio}
\end{figure}

\begin{osservazione}[Dove si trova il minimo]
Un heap binomiale non ha un'unica radice: è una foresta, e ogni albero ha la sua. La proprietà di min-heap vale all'interno di ciascun albero, quindi ogni radice contiene il minimo del proprio albero. Il minimo dell'intero heap è la più piccola delle radici: nella figura~\ref{fig:hb-esempio} le radici sono 10, 1 e 6, e il minimo è 1, la radice del $B_2$.

Il puntatore $\mathit{head}[H]$, invece, non indica il minimo, ma la prima radice della lista, che è ordinata per grado e non per chiave. A differenza di un heap binario, in cui il minimo si trova nell'unica radice, per trovarlo bisogna quindi scorrere tutta la lista delle radici.
\end{osservazione}

\paragraph{Conseguenze.}
Dalla definizione seguono subito due fatti:
\begin{itemize}
    \item In ogni albero binomiale di un heap binomiale, la radice contiene la chiave minima dell'albero. La chiave minima dell'intero heap, quindi, si trova in una delle radici.
    \item Un heap binomiale $H$ con $n$ nodi è formato da al più $\lfloor \lg n \rfloor + 1$ alberi binomiali.
\end{itemize}
Il primo fatto segue dalla proprietà di min-heap: risalendo da un nodo verso la radice le chiavi non aumentano mai. Per dimostrare il secondo, sia $B_k$ l'albero binomiale di grado massimo in $H$. Questo albero ha $2^k$ nodi, che non possono essere più di $n$:
\[
2^k \le n \implies k \le \lg n \implies k \le \lfloor \lg n \rfloor,
\]
dove l'ultimo passaggio vale in quanto $k$ è un intero. Gli alberi di $H$ hanno gradi distinti, compresi tra $0$ e $k$, quindi sono al più $k + 1 \le \lfloor \lg n \rfloor + 1$.

Lo stesso passaggio permette di estendere il corollario~\ref{cor:hb-grado} dal singolo albero binomiale all'intero heap.

\begin{corollario}[Grado massimo in un heap binomiale]\label{cor:hb-grado-heap}
Sia $H$ un heap binomiale con $n \ge 1$ nodi. Allora ogni nodo di $H$ ha grado al più $\lfloor \lg n \rfloor$.
\end{corollario}

\paragraph{Dimostrazione.}
Fissiamo un nodo $x$ di $H$, e sia $B_k$ l'albero binomiale di $H$ che lo contiene. Per le proprietà 4(a) e 4(b) del lemma~\ref{lem:hb-proprieta} ogni nodo di $B_k$ ha grado al più $k$, quindi, qualunque sia la posizione di $x$ in $B_k$:
\[
\deg(x) \le k.
\]
Per la proprietà 1 del lemma~\ref{lem:hb-proprieta}, $B_k$ ha $2^k$ nodi. I nodi di $B_k$ sono anche nodi di $H$, che in totale ne ha $n$, quindi $B_k$ non può averne di più:
\[
2^k \le n.
\]
Applicando $\lg$ a entrambi i lati, che è una funzione crescente, si ottiene
\[
k \le \lg n.
\]
Poiché $k \in \mathbb{N}$ e $\lfloor \lg n \rfloor$ è il più grande intero minore o uguale a $\lg n$, vale anche $k \le \lfloor \lg n \rfloor$. Unendo questa disuguaglianza con la prima:
\[
\deg(x) \le k \le \lfloor \lg n \rfloor.
\]
Poiché $x$ era un nodo qualunque di $H$, la tesi vale per ogni nodo. \hfill $\square$

Rispetto al corollario~\ref{cor:hb-grado} compare la parte intera. In un albero binomiale con $n$ nodi si ha $n = 2^k$, quindi $\lg n$ è già un intero. In un heap binomiale, invece, $n$ è un intero qualsiasi, e l'albero che contiene $x$ ha soltanto $2^k \le n$ nodi.

\begin{osservazione}[La maggiorazione è stretta]
Una maggiorazione si dice \textbf{stretta} se non si può migliorare, ovvero se esiste almeno un caso in cui vale l'uguaglianza. Nel corollario~\ref{cor:hb-grado-heap} l'uguaglianza vale addirittura per ogni heap: un heap binomiale con $n$ nodi contiene sempre un nodo di grado esattamente $\lfloor \lg n \rfloor$.

Sia $B_m$ l'albero binomiale di grado massimo in $H$. Da una parte, come abbiamo visto, $2^m \le n$. Dall'altra, gli alberi di $H$ hanno gradi distinti compresi tra $0$ e $m$. Il numero di nodi, quindi, è massimo quando $H$ li contiene tutti:
\[
n \le 2^0 + 2^1 + \cdots + 2^m = 2^{m+1} - 1 < 2^{m+1}.
\]
Dalle due disuguaglianze si ottiene $2^m \le n < 2^{m+1}$, ovvero $m \le \lg n < m + 1$, e quindi $m = \lfloor \lg n \rfloor$. Per la proprietà 4(a) la radice di $B_m$ ha grado $m$, quindi il grado massimo di un nodo di $H$ è esattamente $\lfloor \lg n \rfloor$.

Ad esempio, l'heap della figura~\ref{fig:hb-esempio} ha $n = 13$ nodi, e $2^3 \le 13 < 2^4$. L'albero di grado massimo è infatti $B_3$, la cui radice $6$ ha tre figli, e $\lfloor \lg 13 \rfloor = 3$. In binario si ha $13 = 1101_2$: gli alberi $B_3$, $B_2$ e $B_0$ dell'heap corrispondono alle cifre uguali a $1$, e l'albero di grado massimo alla cifra più significativa.
\end{osservazione}

\subsection{Rappresentazione}\label{sec:hb-rappresentazione}
Ogni nodo $x$ di un heap binomiale ha i seguenti campi (figura~\ref{fig:hb-rappresentazione}a):
\begin{itemize}
    \item $\mathit{key}[x]$: la chiave, insieme alle eventuali informazioni satellite.
    \item $p[x]$: un puntatore al padre, che vale \Nil{} se $x$ è una radice.
    \item $\mathit{child}[x]$: un puntatore al figlio più a sinistra, che vale \Nil{} se $x$ non ha figli.
    \item $\mathit{sibling}[x]$: un puntatore al fratello immediatamente a destra, che vale \Nil{} se $x$ è l'ultimo figlio di suo padre, oppure l'ultima radice.
    \item $\mathit{degree}[x]$: il grado di $x$, ovvero il numero dei suoi figli.
\end{itemize}
I figli di un nodo formano quindi una lista concatenata tramite il campo $\mathit{sibling}$, e il nodo conosce soltanto il primo.

Anche le radici sono collegate tramite il campo $\mathit{sibling}$, e formano la \textbf{lista delle radici}, ordinata per grado crescente. L'heap è individuato da un puntatore $\mathit{head}[H]$ alla prima radice, che vale \Nil{} se l'heap è vuoto. La figura~\ref{fig:hb-rappresentazione}b mostra la rappresentazione dell'heap di figura~\ref{fig:hb-esempio}.

\begin{figure}[htbp]
\centering
\begin{tikzpicture}[binheap, >=Stealth]
  \node[lab, anchor=west] at (-3.2,1.1) {(a)};
  \draw[Accent, fill=Accent!6] (-1.2,0.96) rectangle (1.2,-0.96);
  \foreach \y in {0.48, 0, -0.48} \draw[Accent] (-1.2,\y) -- (1.2,\y);
  \draw[Accent] (0,-0.48) -- (0,-0.96);
  \node[font=\small] at (0,0.72) {$p$};
  \node[font=\small] at (0,0.24) {$\mathit{key}$};
  \node[font=\small] at (0,-0.24) {$\mathit{degree}$};
  \node[font=\small] at (-0.6,-0.72) {$\mathit{child}$};
  \node[font=\small] at (0.6,-0.72) {$\mathit{sibling}$};
  \draw[->, ptrp] (0,0.96) -- (0,1.45) node[lab, right] {padre};
  \draw[->, ptrc] (-0.6,-0.96) -- (-0.6,-1.45) node[lab, left] {primo figlio};
  \draw[->, ptrs] (1.2,-0.72) -- (1.75,-0.72) node[lab, right] {fratello destro};
\end{tikzpicture}

\vspace{4mm}
\begin{tikzpicture}[binheap, >=Stealth]
  % nodo 10
  \draw[Accent, fill=Accent!6] (-0.400,0.480) rectangle (0.400,-0.480);
  \draw[Accent] (-0.400,0.240) -- (0.400,0.240);
  \draw[Accent] (-0.400,0.000) -- (0.400,0.000);
  \draw[Accent] (-0.400,-0.240) -- (0.400,-0.240);
  \draw[Accent] (0.000,-0.240) -- (0.000,-0.480);
  \node[rec] at (0.000,0.120) {10};
  \node[rec] at (0.000,-0.120) {0};
  \draw[Accent] (-0.340,0.280) -- (0.340,0.440);
  \draw[Accent] (-0.350,-0.440) -- (-0.050,-0.280);
  % nodo 18
  \draw[Accent, fill=Accent!6] (2.085,-2.720) rectangle (2.885,-3.680);
  \draw[Accent] (2.085,-2.960) -- (2.885,-2.960);
  \draw[Accent] (2.085,-3.200) -- (2.885,-3.200);
  \draw[Accent] (2.085,-3.440) -- (2.885,-3.440);
  \draw[Accent] (2.485,-3.440) -- (2.485,-3.680);
  \node[rec] at (2.485,-3.080) {18};
  \node[rec] at (2.485,-3.320) {0};
  \draw[Accent] (2.135,-3.640) -- (2.435,-3.480);
  \draw[Accent] (2.535,-3.640) -- (2.835,-3.480);
  % nodo 12
  \draw[Accent, fill=Accent!6] (2.085,-1.120) rectangle (2.885,-2.080);
  \draw[Accent] (2.085,-1.360) -- (2.885,-1.360);
  \draw[Accent] (2.085,-1.600) -- (2.885,-1.600);
  \draw[Accent] (2.085,-1.840) -- (2.885,-1.840);
  \draw[Accent] (2.485,-1.840) -- (2.485,-2.080);
  \node[rec] at (2.485,-1.480) {12};
  \node[rec] at (2.485,-1.720) {1};
  % nodo 25
  \draw[Accent, fill=Accent!6] (3.505,-1.120) rectangle (4.305,-2.080);
  \draw[Accent] (3.505,-1.360) -- (4.305,-1.360);
  \draw[Accent] (3.505,-1.600) -- (4.305,-1.600);
  \draw[Accent] (3.505,-1.840) -- (4.305,-1.840);
  \draw[Accent] (3.905,-1.840) -- (3.905,-2.080);
  \node[rec] at (3.905,-1.480) {25};
  \node[rec] at (3.905,-1.720) {0};
  \draw[Accent] (3.555,-2.040) -- (3.855,-1.880);
  \draw[Accent] (3.955,-2.040) -- (4.255,-1.880);
  % nodo 1
  \draw[Accent, fill=Accent!6] (3.505,0.480) rectangle (4.305,-0.480);
  \draw[Accent] (3.505,0.240) -- (4.305,0.240);
  \draw[Accent] (3.505,0.000) -- (4.305,0.000);
  \draw[Accent] (3.505,-0.240) -- (4.305,-0.240);
  \draw[Accent] (3.905,-0.240) -- (3.905,-0.480);
  \node[rec] at (3.905,0.120) {1};
  \node[rec] at (3.905,-0.120) {2};
  \draw[Accent] (3.565,0.280) -- (4.245,0.440);
  % nodo 27
  \draw[Accent, fill=Accent!6] (5.990,-4.320) rectangle (6.790,-5.280);
  \draw[Accent] (5.990,-4.560) -- (6.790,-4.560);
  \draw[Accent] (5.990,-4.800) -- (6.790,-4.800);
  \draw[Accent] (5.990,-5.040) -- (6.790,-5.040);
  \draw[Accent] (6.390,-5.040) -- (6.390,-5.280);
  \node[rec] at (6.390,-4.680) {27};
  \node[rec] at (6.390,-4.920) {0};
  \draw[Accent] (6.040,-5.240) -- (6.340,-5.080);
  \draw[Accent] (6.440,-5.240) -- (6.740,-5.080);
  % nodo 11
  \draw[Accent, fill=Accent!6] (5.990,-2.720) rectangle (6.790,-3.680);
  \draw[Accent] (5.990,-2.960) -- (6.790,-2.960);
  \draw[Accent] (5.990,-3.200) -- (6.790,-3.200);
  \draw[Accent] (5.990,-3.440) -- (6.790,-3.440);
  \draw[Accent] (6.390,-3.440) -- (6.390,-3.680);
  \node[rec] at (6.390,-3.080) {11};
  \node[rec] at (6.390,-3.320) {1};
  % nodo 17
  \draw[Accent, fill=Accent!6] (7.410,-2.720) rectangle (8.210,-3.680);
  \draw[Accent] (7.410,-2.960) -- (8.210,-2.960);
  \draw[Accent] (7.410,-3.200) -- (8.210,-3.200);
  \draw[Accent] (7.410,-3.440) -- (8.210,-3.440);
  \draw[Accent] (7.810,-3.440) -- (7.810,-3.680);
  \node[rec] at (7.810,-3.080) {17};
  \node[rec] at (7.810,-3.320) {0};
  \draw[Accent] (7.460,-3.640) -- (7.760,-3.480);
  \draw[Accent] (7.860,-3.640) -- (8.160,-3.480);
  % nodo 8
  \draw[Accent, fill=Accent!6] (7.410,-1.120) rectangle (8.210,-2.080);
  \draw[Accent] (7.410,-1.360) -- (8.210,-1.360);
  \draw[Accent] (7.410,-1.600) -- (8.210,-1.600);
  \draw[Accent] (7.410,-1.840) -- (8.210,-1.840);
  \draw[Accent] (7.810,-1.840) -- (7.810,-2.080);
  \node[rec] at (7.810,-1.480) {8};
  \node[rec] at (7.810,-1.720) {2};
  % nodo 38
  \draw[Accent, fill=Accent!6] (8.830,-2.720) rectangle (9.630,-3.680);
  \draw[Accent] (8.830,-2.960) -- (9.630,-2.960);
  \draw[Accent] (8.830,-3.200) -- (9.630,-3.200);
  \draw[Accent] (8.830,-3.440) -- (9.630,-3.440);
  \draw[Accent] (9.230,-3.440) -- (9.230,-3.680);
  \node[rec] at (9.230,-3.080) {38};
  \node[rec] at (9.230,-3.320) {0};
  \draw[Accent] (8.880,-3.640) -- (9.180,-3.480);
  \draw[Accent] (9.280,-3.640) -- (9.580,-3.480);
  % nodo 14
  \draw[Accent, fill=Accent!6] (8.830,-1.120) rectangle (9.630,-2.080);
  \draw[Accent] (8.830,-1.360) -- (9.630,-1.360);
  \draw[Accent] (8.830,-1.600) -- (9.630,-1.600);
  \draw[Accent] (8.830,-1.840) -- (9.630,-1.840);
  \draw[Accent] (9.230,-1.840) -- (9.230,-2.080);
  \node[rec] at (9.230,-1.480) {14};
  \node[rec] at (9.230,-1.720) {1};
  % nodo 29
  \draw[Accent, fill=Accent!6] (10.250,-1.120) rectangle (11.050,-2.080);
  \draw[Accent] (10.250,-1.360) -- (11.050,-1.360);
  \draw[Accent] (10.250,-1.600) -- (11.050,-1.600);
  \draw[Accent] (10.250,-1.840) -- (11.050,-1.840);
  \draw[Accent] (10.650,-1.840) -- (10.650,-2.080);
  \node[rec] at (10.650,-1.480) {29};
  \node[rec] at (10.650,-1.720) {0};
  \draw[Accent] (10.300,-2.040) -- (10.600,-1.880);
  \draw[Accent] (10.700,-2.040) -- (11.000,-1.880);
  % nodo 6
  \draw[Accent, fill=Accent!6] (10.250,0.480) rectangle (11.050,-0.480);
  \draw[Accent] (10.250,0.240) -- (11.050,0.240);
  \draw[Accent] (10.250,0.000) -- (11.050,0.000);
  \draw[Accent] (10.250,-0.240) -- (11.050,-0.240);
  \draw[Accent] (10.650,-0.240) -- (10.650,-0.480);
  \node[rec] at (10.650,0.120) {6};
  \node[rec] at (10.650,-0.120) {3};
  \draw[Accent] (10.310,0.280) -- (10.990,0.440);
  \draw[Accent] (10.700,-0.440) -- (11.000,-0.280);
  \draw[->, ptrs] (0.200,-0.360) -- (3.505,-0.360);
  \draw[->, ptrp] (2.615,-2.840) -- (2.505,-2.080);
  \draw[->, ptrc] (2.285,-1.960) -- (2.355,-2.720);
  \draw[->, ptrp] (2.615,-1.240) -- (3.925,-0.480);
  \draw[->, ptrs] (2.685,-1.960) -- (3.505,-1.960);
  \draw[->, ptrp] (4.035,-1.240) -- (3.925,-0.480);
  \draw[->, ptrc] (3.705,-0.360) -- (2.355,-1.120);
  \draw[->, ptrs] (4.105,-0.360) -- (10.250,-0.360);
  \draw[->, ptrp] (6.520,-4.440) -- (6.410,-3.680);
  \draw[->, ptrc] (6.190,-3.560) -- (6.260,-4.320);
  \draw[->, ptrp] (6.520,-2.840) -- (7.830,-2.080);
  \draw[->, ptrs] (6.590,-3.560) -- (7.410,-3.560);
  \draw[->, ptrp] (7.940,-2.840) -- (7.830,-2.080);
  \draw[->, ptrc] (7.610,-1.960) -- (6.260,-2.720);
  \draw[->, ptrp] (7.940,-1.240) -- (10.670,-0.480);
  \draw[->, ptrs] (8.010,-1.960) -- (8.830,-1.960);
  \draw[->, ptrp] (9.360,-2.840) -- (9.250,-2.080);
  \draw[->, ptrc] (9.030,-1.960) -- (9.100,-2.720);
  \draw[->, ptrp] (9.360,-1.240) -- (10.670,-0.480);
  \draw[->, ptrs] (9.430,-1.960) -- (10.250,-1.960);
  \draw[->, ptrp] (10.780,-1.240) -- (10.670,-0.480);
  \draw[->, ptrc] (10.450,-0.360) -- (7.680,-1.120);
  \node[lab, anchor=west] at (-2.750,0.950) {(b)};
  \node[head, anchor=east] (rh) at (-1.150,-0.360) {$\mathit{head}[H]$};
  \draw[->, ptrs] (rh) -- (-0.400,-0.360);
\end{tikzpicture}
\caption{(a) I campi di un nodo. (b) La rappresentazione dell'heap di figura~\ref{fig:hb-esempio}: in ogni nodo, dall'alto, il puntatore al padre, la chiave, il grado e i puntatori al primo figlio e al fratello destro. Una barra indica un puntatore che vale \Nil{}. Le radici sono collegate dal campo $\mathit{sibling}$ nella lista delle radici, ordinata per grado crescente, mentre i figli di ogni nodo sono in ordine di grado decrescente.}
\label{fig:hb-rappresentazione}
\end{figure}

\section{Operazioni sugli heap binomiali}\label{sec:hb-op}
Vediamo ora come si realizzano le operazioni dei mergeable heap. Le procedure che studiamo sono le seguenti:
\begin{itemize}
    \item \MakeBinomialHeap$()$: crea un heap binomiale vuoto.
    \item \BinomialHeapMinimum$(H)$: restituisce il nodo con la chiave minima.
    \item \BinomialHeapUnion$(H_1,H_2)$: unisce due heap binomiali, e si appoggia a due procedure ausiliarie, \BinomialLink{} e \BinomialHeapMerge{}.
    \item \BinomialHeapInsert$(H,x)$: inserisce il nodo $x$.
    \item \BinomialHeapExtractMin$(H)$: estrae il nodo con la chiave minima.
    \item \BinomialHeapDecreaseKey$(H,x,k)$: assegna al nodo $x$ la chiave $k$, non maggiore di quella attuale.
    \item \BinomialHeapDelete$(H,x)$: cancella il nodo $x$.
\end{itemize}
In tutti i costi, $n$ indica il numero di nodi dell'heap. Per l'unione, $n$ è la somma dei nodi dei due heap.

\subsection{Creazione di un heap vuoto}\label{sec:hb-make}
Un heap binomiale vuoto è semplicemente un oggetto con $\mathit{head}[H] = \Nil$ (algoritmo~\ref{alg:hb-make}). La procedura richiede tempo $O(1)$.

\begin{algorithm}[htbp]
\caption{La creazione di un heap binomiale vuoto.}\label{alg:hb-make}
\begin{algorithmic}[1]
\Procedure{Make-Binomial-Heap}{}
  \State $H \gets$ \Call{Allocate-Heap}{}
  \State $\mathit{head}[H] \gets \Nil$
  \State \Return $H$
\EndProcedure
\end{algorithmic}
\end{algorithm}

\subsection{Ricerca del minimo}\label{sec:hb-minimum}
Per la proprietà di min-heap, la chiave minima si trova in una delle radici. \BinomialHeapMinimum{} scorre quindi la lista delle radici, e restituisce la radice con la chiave minima (algoritmo~\ref{alg:hb-minimum}). Nell'heap di figura~\ref{fig:hb-esempio} le radici sono 10, 1 e 6, e la procedura restituisce la radice con chiave 1.

\begin{algorithm}[htbp]
\caption{La ricerca del nodo con la chiave minima, che si trova in una delle radici.}\label{alg:hb-minimum}
\begin{algorithmic}[1]
\Procedure{Binomial-Heap-Minimum}{$H$}
  \State $y \gets \Nil$
  \State $x \gets \mathit{head}[H]$
  \State $\mathit{min} \gets \infty$
  \While{$x \ne \Nil$}
    \If{$\mathit{key}[x] < \mathit{min}$}
      \State $\mathit{min} \gets \mathit{key}[x]$
      \State $y \gets x$
    \EndIf
    \State $x \gets \mathit{sibling}[x]$
  \EndWhile
  \State \Return $y$
\EndProcedure
\end{algorithmic}
\end{algorithm}

\paragraph{Costo.}
La procedura esegue un'iterazione per ogni radice, e ogni iterazione costa $O(1)$. Il costo è quindi proporzionale alla lunghezza della lista delle radici, che contiene al più $\lfloor \lg n \rfloor + 1$ alberi:
\[
O(\lfloor \lg n \rfloor + 1) = O(\lg n).
\]

\subsection{Unione}\label{sec:hb-union}
L'unione è l'operazione che distingue gli heap binomiali dagli heap binari, ed è anche quella su cui si appoggiano quasi tutte le altre. La costruiamo in tre passi: il collegamento di due alberi dello stesso grado, la fusione di due liste delle radici, e infine la procedura completa.

\paragraph{Collegare due alberi.}
La procedura \BinomialLink$(y,z)$ riceve le radici $y$ e $z$ di due alberi binomiali dello stesso grado, con $\mathit{key}[z] \le \mathit{key}[y]$, e rende $y$ il primo figlio di $z$ (algoritmo~\ref{alg:hb-link}). È esattamente la costruzione della definizione~\ref{def:albero-binomiale}: due copie di $B_{k-1}$ diventano un $B_k$ (figura~\ref{fig:hb-link}).

\begin{algorithm}[htbp]
\caption{Il collegamento di due alberi binomiali dello stesso grado, con $\mathit{key}[z] \le \mathit{key}[y]$: $y$ diventa il primo figlio di $z$.}\label{alg:hb-link}
\begin{algorithmic}[1]
\Procedure{Binomial-Link}{$y,z$}
  \State $p[y] \gets z$
  \State $\mathit{sibling}[y] \gets \mathit{child}[z]$
  \State $\mathit{child}[z] \gets y$
  \State $\mathit{degree}[z] \gets \mathit{degree}[z] + 1$
\EndProcedure
\end{algorithmic}
\end{algorithm}

Il vecchio primo figlio di $z$ diventa il fratello destro di $y$, e $y$ prende il suo posto in testa alla lista dei figli. Poiché $z$ ha la chiave minore, il nuovo albero rispetta ancora la proprietà di min-heap. La procedura modifica un numero costante di campi, quindi costa $O(1)$.

\begin{figure}[htbp]
\centering
\begin{tikzpicture}[binheap, >=Stealth]
  % \sotto{x}{y}: un sottoalbero B_{k-1} con vertice in (x,y), con il nome sotto la base
  \newcommand{\sotto}[2]{%
    \draw[Accent, fill=Accent!8] (#1,#2) -- ({#1-0.42},{#2-0.95}) -- ({#1+0.42},{#2-0.95}) -- cycle;
    \node[font=\footnotesize] at (#1,{#2-1.17}) {$B_{k-1}$};}
  % prima del collegamento
  \sotto{0}{0}
  \sotto{1.5}{0}
  \node[hn] at (0,0) {$y$};
  \node[hn] at (1.5,0) {$z$};
  \node[lab, font=\footnotesize] at (0.75,-1.7) {$\mathit{key}[z] \le \mathit{key}[y]$};
  \draw[->, thick] (2.5,-0.45) -- node[lab, above] {\BinomialLink$(y,z)$} (4.7,-0.45);
  % dopo il collegamento
  \sotto{5.7}{-0.55}
  \sotto{6.9}{0}
  \node[hl] (y) at (5.7,-0.55) {$y$};
  \node[hl] (z) at (6.9,0) {$z$};
  \draw (z) -- (y);
  \draw[decorate, decoration={brace, mirror, amplitude=4pt}, TextGray] (5.2,-2.05) -- node[lab, below=4pt] {$B_k$} (7.4,-2.05);
\end{tikzpicture}
\caption{\BinomialLink$(y,z)$ collega le radici di due alberi binomiali $B_{k-1}$, con $\mathit{key}[z] \le \mathit{key}[y]$: $y$ diventa il primo figlio di $z$, e si ottiene un $B_k$ con la proprietà di min-heap.}
\label{fig:hb-link}
\end{figure}

\paragraph{L'idea: la somma binaria.}
Per capire l'unione conviene partire da un modo alternativo di sommare due numeri binari. Un albero $B_k$ ha $2^k$ nodi, quindi un heap con $103 = 2^6 + 2^5 + 2^2 + 2^1 + 2^0$ nodi è formato dagli alberi $B_6$, $B_5$, $B_2$, $B_1$ e $B_0$. Sommiamo $n_1 = 103$ e $n_2 = 54$, scrivendo i due numeri come somme di potenze di $2$ e mettendo insieme le potenze in ordine crescente:
\[
\begin{aligned}
103 + 54 &= (2^0 + 2^1 + 2^2 + 2^5 + 2^6) + (2^1 + 2^2 + 2^4 + 2^5) \\
&= 2^0 + 2^1 + 2^1 + 2^2 + 2^2 + 2^4 + 2^5 + 2^5 + 2^6.
\end{aligned}
\]
Questo primo passo corrisponde alla fusione delle due liste delle radici, ordinate per grado. Poi si scorrono le potenze da sinistra a destra e, ogni volta che due potenze uguali sono vicine, le si sostituisce con la loro somma $2^j + 2^j = 2^{j+1}$:
\[
\begin{aligned}
103 + 54 &= 2^0 + 2^2 + 2^2 + 2^2 + 2^4 + 2^5 + 2^5 + 2^6 \\
&= 2^0 + 2^2 + 2^3 + 2^4 + 2^5 + 2^5 + 2^6 \\
&= 2^0 + 2^2 + 2^3 + 2^4 + 2^6 + 2^6 \\
&= 2^0 + 2^2 + 2^3 + 2^4 + 2^7 = 157.
\end{aligned}
\]
Ogni sostituzione $2^j + 2^j = 2^{j+1}$ corrisponde a un \BinomialLink{} tra due alberi $B_j$, che produce un $B_{j+1}$: è il \textbf{riporto} della somma in colonna. Il caso delicato è quello delle tre potenze $2^2$ consecutive, dopo la prima sostituzione. Due vengono dagli heap di partenza e una dal riporto appena creato. La prima resta, e si sommano le altre due, come in una colonna della somma binaria con tre bit uguali a $1$, dove $1 + 1 + 1 = (11)_2$.

La stessa somma, scritta in colonna con i riporti, è la seguente:
\[
\begin{array}{rcccccccc}
 & B_7 & B_6 & B_5 & B_4 & B_3 & B_2 & B_1 & B_0 \\
\text{riporti} & 1 & 1 & & & 1 & 1 & & \\
n_1 = 103 & & 1 & 1 & 0 & 0 & 1 & 1 & 1 \\
n_2 = 54 & & & 1 & 1 & 0 & 1 & 1 & 0 \\
\hline
n = 157 & 1 & 0 & 0 & 1 & 1 & 1 & 0 & 1
\end{array}
\]
dove la colonna $B_i$ contiene l'$i$-esimo bit, ovvero dice se l'heap contiene un albero di grado $i$. I quattro riporti sono i quattro collegamenti che servono per unire i due heap.

\paragraph{Fusione delle liste delle radici.}
La procedura \BinomialHeapMerge$(H_1,H_2)$ fonde le liste delle radici di $H_1$ e di $H_2$ in un'unica lista ordinata per grado crescente, e ne restituisce la prima radice. A parità di grado viene prima la radice di $H_1$. Le due liste sono già ordinate per grado, quindi basta una sola scansione, e il costo è proporzionale al numero complessivo di radici, ovvero $O(\lg n)$.

\paragraph{La procedura completa.}
\BinomialHeapUnion$(H_1,H_2)$ crea un heap vuoto $H$, gli assegna la lista prodotta da \BinomialHeapMerge{}, e poi la scorre da sinistra a destra collegando le radici dello stesso grado, come nella somma vista sopra (algoritmo~\ref{alg:hb-union}).

\begin{algorithm}[htbp]
\caption{L'unione di due heap binomiali: la fusione delle liste delle radici, poi il collegamento delle radici dello stesso grado.}\label{alg:hb-union}
\begin{algorithmic}[1]
\Procedure{Binomial-Heap-Union}{$H_1,H_2$}
  \State $H \gets$ \Call{Make-Binomial-Heap}{}
  \State $\mathit{head}[H] \gets$ \Call{Binomial-Heap-Merge}{$H_1,H_2$}
  \State libera gli oggetti $H_1$ e $H_2$, ma non le liste a cui puntano
  \If{$\mathit{head}[H] = \Nil$}
    \State \Return $H$
  \EndIf
  \State $\prevx \gets \Nil$
  \State $x \gets \mathit{head}[H]$
  \State $\nextx \gets \mathit{sibling}[x]$
  \While{$\nextx \ne \Nil$}
    \If{$\mathit{degree}[x] \ne \mathit{degree}[\nextx]$}
      \State $\prevx \gets x$ \Comment{caso 1}
      \State $x \gets \nextx$
    \ElsIf{$\mathit{sibling}[\nextx] \ne \Nil$ \textbf{and} $\mathit{degree}[\mathit{sibling}[\nextx]] = \mathit{degree}[x]$}
      \State $\prevx \gets x$ \Comment{caso 2}
      \State $x \gets \nextx$
    \ElsIf{$\mathit{key}[x] \le \mathit{key}[\nextx]$}
      \State $\mathit{sibling}[x] \gets \mathit{sibling}[\nextx]$ \Comment{caso 3}
      \State \Call{Binomial-Link}{$\nextx, x$}
    \Else
      \If{$\prevx = \Nil$} \Comment{caso 4}
        \State $\mathit{head}[H] \gets \nextx$
      \Else
        \State $\mathit{sibling}[\prevx] \gets \nextx$
      \EndIf
      \State \Call{Binomial-Link}{$x, \nextx$}
      \State $x \gets \nextx$
    \EndIf
    \State $\nextx \gets \mathit{sibling}[x]$
  \EndWhile
  \State \Return $H$
\EndProcedure
\end{algorithmic}
\end{algorithm}

Il ciclo mantiene tre puntatori a radici consecutive: $x$ è la radice corrente, $\nextx$ la successiva e $\prevx$ la precedente, che serve per aggiornare la lista quando si toglie $x$. Nella lista fusa ci sono al più due radici di ogni grado. Collegando due radici di grado $k$, però, se ne ottiene una di grado $k+1$, che può trovarsi accanto alle due radici di grado $k+1$ dei due heap. In ogni momento, quindi, ci sono al più tre radici consecutive dello stesso grado.

A ogni iterazione si guardano i gradi di $x$, di $\nextx$ e di $\mathit{sibling}[\nextx]$, e si distinguono quattro casi.

\begin{description}
  \item[{Caso 1: $\mathit{degree}[x] \ne \mathit{degree}[\nextx]$.}] Sia $x$ la radice di un $B_k$ e $\nextx$ la radice di un $B_l$, con $l > k$. Allora $x$ è l'unica radice di grado $k$, e non c'è niente da collegare: si avanza di una posizione, e $x$ passa alla radice successiva. In binario è una colonna con un solo bit uguale a $1$, che resta nel risultato.

  \begin{casofig}[1]
  I puntatori avanzano di una posizione, e la lista non cambia. Lo stesso accade quando $\nextx$ è l'ultima radice della lista: dopo l'avanzamento $\nextx$ diventa \Nil{}, e il ciclo termina.
  \tcblower
  \begin{widget}{union-case}[case=1]
  \begin{tikzpicture}[binheap, >=Stealth]
    % \rad{nome}{etichetta}{x}{y}{puntatore}{albero}: una radice con il suo sottoalbero
    \newcommand{\rad}[6]{%
      \draw[Accent, fill=Accent!6] (#3,#4) -- ({#3-0.24},{#4-0.62}) -- ({#3+0.24},{#4-0.62}) -- cycle;
      \node[hn] (#1) at (#3,#4) {$#2$};
      \node[ptr, anchor=south] at (#1.north) {#5};
      \node[font=\scriptsize] at (#3,{#4-0.82}) {$#6$};}
    \rad{a1}{a}{0}{0}{\textit{prev-x}}{}
    \rad{b1}{b}{1.4}{0}{\textit{x}}{B_k}
    \rad{c1}{c}{2.8}{0}{\textit{next-x}}{B_l}
    \rad{d1}{d}{4.2}{0}{\textit{sibling}[\textit{next-x}]}{}
    \draw[->] (a1) -- (b1); \draw[->] (b1) -- (c1); \draw[->] (c1) -- (d1);
    \draw[->, thick] (2.1,-1.15) -- node[lab, right, font=\scriptsize] {caso 1} (2.1,-1.6);
    \rad{a2}{a}{0}{-2.2}{}{}
    \rad{b2}{b}{1.4}{-2.2}{\textit{prev-x}}{B_k}
    \rad{c2}{c}{2.8}{-2.2}{\textit{x}}{B_l}
    \rad{d2}{d}{4.2}{-2.2}{\textit{next-x}}{}
    \draw[->] (a2) -- (b2); \draw[->] (b2) -- (c2); \draw[->] (c2) -- (d2);
  \end{tikzpicture}
  \end{widget}
  \end{casofig}

  \item[{Caso 2: tre radici consecutive dello stesso grado.}] Le radici $x$, $\nextx$ e $\mathit{sibling}[\nextx]$ hanno tutte grado $k$: la prima è il riporto prodotto dall'ultimo collegamento, le altre due vengono dai due heap. Si avanza di una posizione, e alla prossima iterazione si collegano le ultime due, con il caso 3 o 4. In binario è la colonna con tre bit uguali a $1$: un bit resta nel risultato, e uno va in riporto.

  \begin{casofig}[2]
  Come nel caso 1, i puntatori avanzano di una posizione e la lista non cambia. La radice $b$ resta nel risultato come unico $B_k$, mentre $c$ e $d$ verranno collegate.
  \tcblower
  \begin{widget}{union-case}[case=2]
  \begin{tikzpicture}[binheap, >=Stealth]
    \newcommand{\rad}[6]{%
      \draw[Accent, fill=Accent!6] (#3,#4) -- ({#3-0.24},{#4-0.62}) -- ({#3+0.24},{#4-0.62}) -- cycle;
      \node[hn] (#1) at (#3,#4) {$#2$};
      \node[ptr, anchor=south] at (#1.north) {#5};
      \node[font=\scriptsize] at (#3,{#4-0.82}) {$#6$};}
    \rad{a1}{a}{0}{0}{\textit{prev-x}}{}
    \rad{b1}{b}{1.4}{0}{\textit{x}}{B_k}
    \rad{c1}{c}{2.8}{0}{\textit{next-x}}{B_k}
    \rad{d1}{d}{4.2}{0}{\textit{sibling}[\textit{next-x}]}{B_k}
    \draw[->] (a1) -- (b1); \draw[->] (b1) -- (c1); \draw[->] (c1) -- (d1);
    \draw[->, thick] (2.1,-1.15) -- node[lab, right, font=\scriptsize] {caso 2} (2.1,-1.6);
    \rad{a2}{a}{0}{-2.2}{}{}
    \rad{b2}{b}{1.4}{-2.2}{\textit{prev-x}}{B_k}
    \rad{c2}{c}{2.8}{-2.2}{\textit{x}}{B_k}
    \rad{d2}{d}{4.2}{-2.2}{\textit{next-x}}{B_k}
    \draw[->] (a2) -- (b2); \draw[->] (b2) -- (c2); \draw[->] (c2) -- (d2);
  \end{tikzpicture}
  \end{widget}
  \end{casofig}

  \item[{Caso 3: due radici dello stesso grado, con $\mathit{key}[x] \le \mathit{key}[\nextx]$.}] Le radici $x$ e $\nextx$ hanno grado $k$, mentre $\mathit{sibling}[\nextx]$ ha un grado diverso oppure è \Nil{}. Sono quindi le uniche due radici di grado $k$, e vanno collegate. Poiché $x$ ha la chiave minore, $\nextx$ viene tolta dalla lista e diventa il primo figlio di $x$, che ora è la radice di un $B_{k+1}$.

  \begin{casofig}[3]
  Si pone $\mathit{sibling}[x] \gets \mathit{sibling}[\nextx]$, che toglie $c$ dalla lista, e si chiama \BinomialLink$(\nextx, x)$. La radice corrente resta $x$: alla prossima iterazione il nuovo $B_{k+1}$ si confronta con la radice successiva, che potrebbe avere anch'essa grado $k+1$.
  \tcblower
  \begin{widget}{union-case}[case=3]
  \begin{tikzpicture}[binheap, >=Stealth]
    \newcommand{\rad}[6]{%
      \draw[Accent, fill=Accent!6] (#3,#4) -- ({#3-0.24},{#4-0.62}) -- ({#3+0.24},{#4-0.62}) -- cycle;
      \node[hn] (#1) at (#3,#4) {$#2$};
      \node[ptr, anchor=south] at (#1.north) {#5};
      \node[font=\scriptsize] at (#3,{#4-0.82}) {$#6$};}
    \rad{a1}{a}{0}{0}{\textit{prev-x}}{}
    \rad{b1}{b}{1.4}{0}{\textit{x}}{B_k}
    \rad{c1}{c}{2.8}{0}{\textit{next-x}}{B_k}
    \rad{d1}{d}{4.2}{0}{\textit{sibling}[\textit{next-x}]}{B_l}
    \draw[->] (a1) -- (b1); \draw[->] (b1) -- (c1); \draw[->] (c1) -- (d1);
    \node[font=\scriptsize] at (2.1,-1.15) {$\mathit{key}[x] \le \mathit{key}[\textit{next-x}]$};
    \draw[->, thick] (2.1,-1.4) -- node[lab, right, font=\scriptsize] {caso 3} (2.1,-1.85);
    \rad{a2}{a}{0}{-2.45}{\textit{prev-x}}{}
    \rad{c2}{c}{1.05}{-3.0}{}{B_k}
    \rad{b2}{b}{1.6}{-2.45}{\textit{x}}{B_k}
    \draw (b2) -- (c2);
    \rad{d2}{d}{3.4}{-2.45}{\textit{next-x}}{B_l}
    \draw[->] (a2) -- (b2); \draw[->] (b2) -- (d2);
    \draw[decorate, decoration={brace, mirror, amplitude=3pt}, TextGray] (0.75,-4.05) -- node[font=\scriptsize, below=3pt] {$B_{k+1}$} (1.9,-4.05);
  \end{tikzpicture}
  \end{widget}
  \end{casofig}

  \item[{Caso 4: due radici dello stesso grado, con $\mathit{key}[x] > \mathit{key}[\nextx]$.}] Questa volta è $\nextx$ ad avere la chiave minore, quindi $x$ viene tolta dalla lista e diventa il primo figlio di $\nextx$. Se $x$ era la prima radice si aggiorna $\mathit{head}[H]$, altrimenti il fratello destro di $\prevx$.

  \begin{casofig}[4]
  Si toglie $x$ dalla lista con $\mathit{head}[H] \gets \nextx$ oppure con $\mathit{sibling}[\prevx] \gets \nextx$, si chiama \BinomialLink$(x, \nextx)$, e infine $x$ avanza sulla radice del nuovo $B_{k+1}$.
  \tcblower
  \begin{widget}{union-case}[case=4]
  \begin{tikzpicture}[binheap, >=Stealth]
    \newcommand{\rad}[6]{%
      \draw[Accent, fill=Accent!6] (#3,#4) -- ({#3-0.24},{#4-0.62}) -- ({#3+0.24},{#4-0.62}) -- cycle;
      \node[hn] (#1) at (#3,#4) {$#2$};
      \node[ptr, anchor=south] at (#1.north) {#5};
      \node[font=\scriptsize] at (#3,{#4-0.82}) {$#6$};}
    \rad{a1}{a}{0}{0}{\textit{prev-x}}{}
    \rad{b1}{b}{1.4}{0}{\textit{x}}{B_k}
    \rad{c1}{c}{2.8}{0}{\textit{next-x}}{B_k}
    \rad{d1}{d}{4.2}{0}{\textit{sibling}[\textit{next-x}]}{B_l}
    \draw[->] (a1) -- (b1); \draw[->] (b1) -- (c1); \draw[->] (c1) -- (d1);
    \node[font=\scriptsize] at (2.1,-1.15) {$\mathit{key}[x] > \mathit{key}[\textit{next-x}]$};
    \draw[->, thick] (2.1,-1.4) -- node[lab, right, font=\scriptsize] {caso 4} (2.1,-1.85);
    \rad{a2}{a}{0}{-2.45}{\textit{prev-x}}{}
    \rad{b2}{b}{1.05}{-3.0}{}{B_k}
    \rad{c2}{c}{1.6}{-2.45}{\textit{x}}{B_k}
    \draw (c2) -- (b2);
    \rad{d2}{d}{3.4}{-2.45}{\textit{next-x}}{B_l}
    \draw[->] (a2) -- (c2); \draw[->] (c2) -- (d2);
    \draw[decorate, decoration={brace, mirror, amplitude=3pt}, TextGray] (0.75,-4.05) -- node[font=\scriptsize, below=3pt] {$B_{k+1}$} (1.9,-4.05);
  \end{tikzpicture}
  \end{widget}
  \end{casofig}
\end{description}

In tutti i casi, alla fine dell'iterazione $\nextx$ diventa la radice che segue $x$. Nei casi 3 e 4, inoltre, $x$ è la radice di un $B_{k+1}$ appena formato: se anche la radice successiva ha grado $k+1$, la prossima iterazione la collegherà a $x$. È il riporto che si propaga verso i gradi più alti.

\paragraph{Esempio.}
Uniamo i due heap della figura~\ref{fig:hb-unione}a, con $n_1 = 7$ e $n_2 = 19$ nodi. \BinomialHeapMerge{} produce la lista delle radici 12, 18, 7, 3, 15 e 6, di grado $0, 0, 1, 1, 2, 4$ (figura~\ref{fig:hb-unione}b). Poi il ciclo esegue i seguenti passi:
\begin{itemize}
    \item \textbf{Caso 3}: 12 e 18 hanno grado $0$, mentre la radice successiva, 7, ha grado $1$. Poiché $12 \le 18$, 18 diventa figlio di 12, che è ora la radice di un $B_1$ (c).
    \item \textbf{Caso 2}: 12, 7 e 3 hanno tutti grado $1$, quindi $x$ avanza su 7. Il $B_1$ con radice 12 resterà nel risultato (d).
    \item \textbf{Caso 4}: 7 e 3 hanno grado $1$, e la radice successiva, 15, ha grado $2$. Poiché $7 > 3$, 7 diventa figlio di 3, e $x$ avanza su 3, ora la radice di un $B_2$ (e).
    \item \textbf{Caso 3}: 3 e 15 hanno grado $2$, mentre 6 ha grado $4$. Poiché $3 \le 15$, 15 diventa figlio di 3, ora la radice di un $B_3$ (f).
    \item \textbf{Caso 1}: 3 ha grado $3$ e 6 ha grado $4$, quindi $x$ avanza su 6. Ora $\nextx = \Nil$, e il ciclo termina.
\end{itemize}
L'heap risultante ha $26$ nodi, e contiene gli alberi $B_1$, $B_3$ e $B_4$. In binario, infatti:
\[
n_1 + n_2 = (00111)_2 + (10011)_2 = (11010)_2 = 26.
\]

\begin{figure}[htbp]
\centering
\begin{widget}{binomial-heap}[heap={(12) (7 (25)) (15 (28 (41)) (33))}, other={(18) (3 (37)) (6 (8 (30 (45 (55)) (32)) (23 (24)) (22)) (29 (48 (50)) (31)) (10 (17)) (44))}, ops={union}, actions={}]
\begin{tikzpicture}[binheap, >=Stealth]
  % (a) 
  \node[lab, anchor=west] at (-1.05,0.000) {(a)};
  \node[hn] (s0n0) at (1.530,0.000) {12};
  \node[hn] (s0n2) at (2.510,-0.600) {25};
  \node[hn] (s0n1) at (2.510,0.000) {7};
  \node[hn] (s0n5) at (3.490,-1.200) {41};
  \node[hn] (s0n4) at (3.490,-0.600) {28};
  \node[hn] (s0n6) at (4.050,-0.600) {33};
  \node[hn] (s0n3) at (4.050,0.000) {15};
  \node[hn] (s0n7) at (6.228,0.000) {18};
  \node[hn] (s0n9) at (7.208,-0.600) {37};
  \node[hn] (s0n8) at (7.208,0.000) {3};
  \node[hn] (s0n14) at (8.188,-2.400) {55};
  \node[hn] (s0n13) at (8.188,-1.800) {45};
  \node[hn] (s0n15) at (8.748,-1.800) {32};
  \node[hn] (s0n12) at (8.748,-1.200) {30};
  \node[hn] (s0n17) at (9.308,-1.800) {24};
  \node[hn] (s0n16) at (9.308,-1.200) {23};
  \node[hn] (s0n18) at (9.868,-1.200) {22};
  \node[hn] (s0n11) at (9.868,-0.600) {8};
  \node[hn] (s0n21) at (10.428,-1.800) {50};
  \node[hn] (s0n20) at (10.428,-1.200) {48};
  \node[hn] (s0n22) at (10.988,-1.200) {31};
  \node[hn] (s0n19) at (10.988,-0.600) {29};
  \node[hn] (s0n24) at (11.548,-1.200) {17};
  \node[hn] (s0n23) at (11.548,-0.600) {10};
  \node[hn] (s0n25) at (12.108,-0.600) {44};
  \node[hn] (s0n10) at (12.108,0.000) {6};
  \draw (s0n1) -- (s0n2);
  \draw (s0n4) -- (s0n5);
  \draw (s0n3) -- (s0n4);
  \draw (s0n3) -- (s0n6);
  \draw (s0n8) -- (s0n9);
  \draw (s0n13) -- (s0n14);
  \draw (s0n12) -- (s0n13);
  \draw (s0n12) -- (s0n15);
  \draw (s0n11) -- (s0n12);
  \draw (s0n16) -- (s0n17);
  \draw (s0n11) -- (s0n16);
  \draw (s0n11) -- (s0n18);
  \draw (s0n10) -- (s0n11);
  \draw (s0n20) -- (s0n21);
  \draw (s0n19) -- (s0n20);
  \draw (s0n19) -- (s0n22);
  \draw (s0n10) -- (s0n19);
  \draw (s0n23) -- (s0n24);
  \draw (s0n10) -- (s0n23);
  \draw (s0n10) -- (s0n25);
  \draw[->] (s0n0) -- (s0n1);
  \draw[->] (s0n1) -- (s0n3);
  \draw[->] (s0n7) -- (s0n8);
  \draw[->] (s0n8) -- (s0n10);
  \node[head, anchor=east] (s0h0) at (1.094,0.000) {$\mathit{head}[H_1]$};
  \draw[->] (s0h0) -- (s0n0);
  \node[head, anchor=east] (s0h1) at (5.792,0.000) {$\mathit{head}[H_2]$};
  \draw[->] (s0h1) -- (s0n7);
  % (b) 
  \node[lab, anchor=west] at (-1.05,-3.250) {(b)};
  \node[hn] (s1n0) at (1.530,-3.250) {12};
  \node[hn] (s1n7) at (2.510,-3.250) {18};
  \node[hn] (s1n2) at (3.490,-3.850) {25};
  \node[hn] (s1n1) at (3.490,-3.250) {7};
  \node[hn] (s1n9) at (4.470,-3.850) {37};
  \node[hn] (s1n8) at (4.470,-3.250) {3};
  \node[hn] (s1n5) at (5.450,-4.450) {41};
  \node[hn] (s1n4) at (5.450,-3.850) {28};
  \node[hn] (s1n6) at (6.010,-3.850) {33};
  \node[hn] (s1n3) at (6.010,-3.250) {15};
  \node[hn] (s1n14) at (6.990,-5.650) {55};
  \node[hn] (s1n13) at (6.990,-5.050) {45};
  \node[hn] (s1n15) at (7.550,-5.050) {32};
  \node[hn] (s1n12) at (7.550,-4.450) {30};
  \node[hn] (s1n17) at (8.110,-5.050) {24};
  \node[hn] (s1n16) at (8.110,-4.450) {23};
  \node[hn] (s1n18) at (8.670,-4.450) {22};
  \node[hn] (s1n11) at (8.670,-3.850) {8};
  \node[hn] (s1n21) at (9.230,-5.050) {50};
  \node[hn] (s1n20) at (9.230,-4.450) {48};
  \node[hn] (s1n22) at (9.790,-4.450) {31};
  \node[hn] (s1n19) at (9.790,-3.850) {29};
  \node[hn] (s1n24) at (10.350,-4.450) {17};
  \node[hn] (s1n23) at (10.350,-3.850) {10};
  \node[hn] (s1n25) at (10.910,-3.850) {44};
  \node[hn] (s1n10) at (10.910,-3.250) {6};
  \draw (s1n1) -- (s1n2);
  \draw (s1n8) -- (s1n9);
  \draw (s1n4) -- (s1n5);
  \draw (s1n3) -- (s1n4);
  \draw (s1n3) -- (s1n6);
  \draw (s1n13) -- (s1n14);
  \draw (s1n12) -- (s1n13);
  \draw (s1n12) -- (s1n15);
  \draw (s1n11) -- (s1n12);
  \draw (s1n16) -- (s1n17);
  \draw (s1n11) -- (s1n16);
  \draw (s1n11) -- (s1n18);
  \draw (s1n10) -- (s1n11);
  \draw (s1n20) -- (s1n21);
  \draw (s1n19) -- (s1n20);
  \draw (s1n19) -- (s1n22);
  \draw (s1n10) -- (s1n19);
  \draw (s1n23) -- (s1n24);
  \draw (s1n10) -- (s1n23);
  \draw (s1n10) -- (s1n25);
  \draw[->] (s1n0) -- (s1n7);
  \draw[->] (s1n7) -- (s1n1);
  \draw[->] (s1n1) -- (s1n8);
  \draw[->] (s1n8) -- (s1n3);
  \draw[->] (s1n3) -- (s1n10);
  \node[head, anchor=east] (s1h0) at (1.094,-3.250) {$\mathit{head}[H]$};
  \draw[->] (s1h0) -- (s1n0);
  \node[ptr, anchor=south] at (s1n0.north) {\textit{x}};
  \node[ptr, anchor=south] at (s1n7.north) {\textit{next-x}};
  % (c) 
  \node[lab, anchor=west] at (-1.05,-6.500) {(c)};
  \node[hl] (s2n7) at (1.530,-7.100) {18};
  \node[hl] (s2n0) at (1.530,-6.500) {12};
  \node[hn] (s2n2) at (2.510,-7.100) {25};
  \node[hn] (s2n1) at (2.510,-6.500) {7};
  \node[hn] (s2n9) at (3.490,-7.100) {37};
  \node[hn] (s2n8) at (3.490,-6.500) {3};
  \node[hn] (s2n5) at (4.470,-7.700) {41};
  \node[hn] (s2n4) at (4.470,-7.100) {28};
  \node[hn] (s2n6) at (5.030,-7.100) {33};
  \node[hn] (s2n3) at (5.030,-6.500) {15};
  \node[hn] (s2n14) at (6.010,-8.900) {55};
  \node[hn] (s2n13) at (6.010,-8.300) {45};
  \node[hn] (s2n15) at (6.570,-8.300) {32};
  \node[hn] (s2n12) at (6.570,-7.700) {30};
  \node[hn] (s2n17) at (7.130,-8.300) {24};
  \node[hn] (s2n16) at (7.130,-7.700) {23};
  \node[hn] (s2n18) at (7.690,-7.700) {22};
  \node[hn] (s2n11) at (7.690,-7.100) {8};
  \node[hn] (s2n21) at (8.250,-8.300) {50};
  \node[hn] (s2n20) at (8.250,-7.700) {48};
  \node[hn] (s2n22) at (8.810,-7.700) {31};
  \node[hn] (s2n19) at (8.810,-7.100) {29};
  \node[hn] (s2n24) at (9.370,-7.700) {17};
  \node[hn] (s2n23) at (9.370,-7.100) {10};
  \node[hn] (s2n25) at (9.930,-7.100) {44};
  \node[hn] (s2n10) at (9.930,-6.500) {6};
  \draw (s2n0) -- (s2n7);
  \draw (s2n1) -- (s2n2);
  \draw (s2n8) -- (s2n9);
  \draw (s2n4) -- (s2n5);
  \draw (s2n3) -- (s2n4);
  \draw (s2n3) -- (s2n6);
  \draw (s2n13) -- (s2n14);
  \draw (s2n12) -- (s2n13);
  \draw (s2n12) -- (s2n15);
  \draw (s2n11) -- (s2n12);
  \draw (s2n16) -- (s2n17);
  \draw (s2n11) -- (s2n16);
  \draw (s2n11) -- (s2n18);
  \draw (s2n10) -- (s2n11);
  \draw (s2n20) -- (s2n21);
  \draw (s2n19) -- (s2n20);
  \draw (s2n19) -- (s2n22);
  \draw (s2n10) -- (s2n19);
  \draw (s2n23) -- (s2n24);
  \draw (s2n10) -- (s2n23);
  \draw (s2n10) -- (s2n25);
  \draw[->] (s2n0) -- (s2n1);
  \draw[->] (s2n1) -- (s2n8);
  \draw[->] (s2n8) -- (s2n3);
  \draw[->] (s2n3) -- (s2n10);
  \node[head, anchor=east] (s2h0) at (1.094,-6.500) {$\mathit{head}[H]$};
  \draw[->] (s2h0) -- (s2n0);
  \node[ptr, anchor=south] at (s2n0.north) {\textit{x}};
  \node[ptr, anchor=south] at (s2n1.north) {\textit{next-x}};
  % (d) 
  \node[lab, anchor=west] at (-1.05,-9.750) {(d)};
  \node[hn] (s3n7) at (1.530,-10.350) {18};
  \node[hn] (s3n0) at (1.530,-9.750) {12};
  \node[hn] (s3n2) at (2.510,-10.350) {25};
  \node[hn] (s3n1) at (2.510,-9.750) {7};
  \node[hn] (s3n9) at (3.490,-10.350) {37};
  \node[hn] (s3n8) at (3.490,-9.750) {3};
  \node[hn] (s3n5) at (4.470,-10.950) {41};
  \node[hn] (s3n4) at (4.470,-10.350) {28};
  \node[hn] (s3n6) at (5.030,-10.350) {33};
  \node[hn] (s3n3) at (5.030,-9.750) {15};
  \node[hn] (s3n14) at (6.010,-12.150) {55};
  \node[hn] (s3n13) at (6.010,-11.550) {45};
  \node[hn] (s3n15) at (6.570,-11.550) {32};
  \node[hn] (s3n12) at (6.570,-10.950) {30};
  \node[hn] (s3n17) at (7.130,-11.550) {24};
  \node[hn] (s3n16) at (7.130,-10.950) {23};
  \node[hn] (s3n18) at (7.690,-10.950) {22};
  \node[hn] (s3n11) at (7.690,-10.350) {8};
  \node[hn] (s3n21) at (8.250,-11.550) {50};
  \node[hn] (s3n20) at (8.250,-10.950) {48};
  \node[hn] (s3n22) at (8.810,-10.950) {31};
  \node[hn] (s3n19) at (8.810,-10.350) {29};
  \node[hn] (s3n24) at (9.370,-10.950) {17};
  \node[hn] (s3n23) at (9.370,-10.350) {10};
  \node[hn] (s3n25) at (9.930,-10.350) {44};
  \node[hn] (s3n10) at (9.930,-9.750) {6};
  \draw (s3n0) -- (s3n7);
  \draw (s3n1) -- (s3n2);
  \draw (s3n8) -- (s3n9);
  \draw (s3n4) -- (s3n5);
  \draw (s3n3) -- (s3n4);
  \draw (s3n3) -- (s3n6);
  \draw (s3n13) -- (s3n14);
  \draw (s3n12) -- (s3n13);
  \draw (s3n12) -- (s3n15);
  \draw (s3n11) -- (s3n12);
  \draw (s3n16) -- (s3n17);
  \draw (s3n11) -- (s3n16);
  \draw (s3n11) -- (s3n18);
  \draw (s3n10) -- (s3n11);
  \draw (s3n20) -- (s3n21);
  \draw (s3n19) -- (s3n20);
  \draw (s3n19) -- (s3n22);
  \draw (s3n10) -- (s3n19);
  \draw (s3n23) -- (s3n24);
  \draw (s3n10) -- (s3n23);
  \draw (s3n10) -- (s3n25);
  \draw[->] (s3n0) -- (s3n1);
  \draw[->] (s3n1) -- (s3n8);
  \draw[->] (s3n8) -- (s3n3);
  \draw[->] (s3n3) -- (s3n10);
  \node[head, anchor=east] (s3h0) at (1.094,-9.750) {$\mathit{head}[H]$};
  \draw[->] (s3h0) -- (s3n0);
  \node[ptr, anchor=south] at (s3n0.north) {\textit{prev-x}};
  \node[ptr, anchor=south] at (s3n1.north) {\textit{x}};
  \node[ptr, anchor=south] at (s3n8.north) {\textit{next-x}};
  % (e) 
  \node[lab, anchor=west] at (-1.05,-13.000) {(e)};
  \node[hn] (s4n7) at (1.530,-13.600) {18};
  \node[hn] (s4n0) at (1.530,-13.000) {12};
  \node[hn] (s4n2) at (2.510,-14.200) {25};
  \node[hl] (s4n1) at (2.510,-13.600) {7};
  \node[hn] (s4n9) at (3.070,-13.600) {37};
  \node[hl] (s4n8) at (3.070,-13.000) {3};
  \node[hn] (s4n5) at (4.050,-14.200) {41};
  \node[hn] (s4n4) at (4.050,-13.600) {28};
  \node[hn] (s4n6) at (4.610,-13.600) {33};
  \node[hn] (s4n3) at (4.610,-13.000) {15};
  \node[hn] (s4n14) at (5.590,-15.400) {55};
  \node[hn] (s4n13) at (5.590,-14.800) {45};
  \node[hn] (s4n15) at (6.150,-14.800) {32};
  \node[hn] (s4n12) at (6.150,-14.200) {30};
  \node[hn] (s4n17) at (6.710,-14.800) {24};
  \node[hn] (s4n16) at (6.710,-14.200) {23};
  \node[hn] (s4n18) at (7.270,-14.200) {22};
  \node[hn] (s4n11) at (7.270,-13.600) {8};
  \node[hn] (s4n21) at (7.830,-14.800) {50};
  \node[hn] (s4n20) at (7.830,-14.200) {48};
  \node[hn] (s4n22) at (8.390,-14.200) {31};
  \node[hn] (s4n19) at (8.390,-13.600) {29};
  \node[hn] (s4n24) at (8.950,-14.200) {17};
  \node[hn] (s4n23) at (8.950,-13.600) {10};
  \node[hn] (s4n25) at (9.510,-13.600) {44};
  \node[hn] (s4n10) at (9.510,-13.000) {6};
  \draw (s4n0) -- (s4n7);
  \draw (s4n1) -- (s4n2);
  \draw (s4n8) -- (s4n1);
  \draw (s4n8) -- (s4n9);
  \draw (s4n4) -- (s4n5);
  \draw (s4n3) -- (s4n4);
  \draw (s4n3) -- (s4n6);
  \draw (s4n13) -- (s4n14);
  \draw (s4n12) -- (s4n13);
  \draw (s4n12) -- (s4n15);
  \draw (s4n11) -- (s4n12);
  \draw (s4n16) -- (s4n17);
  \draw (s4n11) -- (s4n16);
  \draw (s4n11) -- (s4n18);
  \draw (s4n10) -- (s4n11);
  \draw (s4n20) -- (s4n21);
  \draw (s4n19) -- (s4n20);
  \draw (s4n19) -- (s4n22);
  \draw (s4n10) -- (s4n19);
  \draw (s4n23) -- (s4n24);
  \draw (s4n10) -- (s4n23);
  \draw (s4n10) -- (s4n25);
  \draw[->] (s4n0) -- (s4n8);
  \draw[->] (s4n8) -- (s4n3);
  \draw[->] (s4n3) -- (s4n10);
  \node[head, anchor=east] (s4h0) at (1.094,-13.000) {$\mathit{head}[H]$};
  \draw[->] (s4h0) -- (s4n0);
  \node[ptr, anchor=south] at (s4n0.north) {\textit{prev-x}};
  \node[ptr, anchor=south] at (s4n8.north) {\textit{x}};
  \node[ptr, anchor=south] at (s4n3.north) {\textit{next-x}};
  % (f) 
  \node[lab, anchor=west] at (-1.05,-16.250) {(f)};
  \node[hn] (s5n7) at (1.530,-16.850) {18};
  \node[hn] (s5n0) at (1.530,-16.250) {12};
  \node[hn] (s5n5) at (2.510,-18.050) {41};
  \node[hn] (s5n4) at (2.510,-17.450) {28};
  \node[hn] (s5n6) at (3.070,-17.450) {33};
  \node[hl] (s5n3) at (3.070,-16.850) {15};
  \node[hn] (s5n2) at (3.630,-17.450) {25};
  \node[hn] (s5n1) at (3.630,-16.850) {7};
  \node[hn] (s5n9) at (4.190,-16.850) {37};
  \node[hl] (s5n8) at (4.190,-16.250) {3};
  \node[hn] (s5n14) at (5.170,-18.650) {55};
  \node[hn] (s5n13) at (5.170,-18.050) {45};
  \node[hn] (s5n15) at (5.730,-18.050) {32};
  \node[hn] (s5n12) at (5.730,-17.450) {30};
  \node[hn] (s5n17) at (6.290,-18.050) {24};
  \node[hn] (s5n16) at (6.290,-17.450) {23};
  \node[hn] (s5n18) at (6.850,-17.450) {22};
  \node[hn] (s5n11) at (6.850,-16.850) {8};
  \node[hn] (s5n21) at (7.410,-18.050) {50};
  \node[hn] (s5n20) at (7.410,-17.450) {48};
  \node[hn] (s5n22) at (7.970,-17.450) {31};
  \node[hn] (s5n19) at (7.970,-16.850) {29};
  \node[hn] (s5n24) at (8.530,-17.450) {17};
  \node[hn] (s5n23) at (8.530,-16.850) {10};
  \node[hn] (s5n25) at (9.090,-16.850) {44};
  \node[hn] (s5n10) at (9.090,-16.250) {6};
  \draw (s5n0) -- (s5n7);
  \draw (s5n4) -- (s5n5);
  \draw (s5n3) -- (s5n4);
  \draw (s5n3) -- (s5n6);
  \draw (s5n8) -- (s5n3);
  \draw (s5n1) -- (s5n2);
  \draw (s5n8) -- (s5n1);
  \draw (s5n8) -- (s5n9);
  \draw (s5n13) -- (s5n14);
  \draw (s5n12) -- (s5n13);
  \draw (s5n12) -- (s5n15);
  \draw (s5n11) -- (s5n12);
  \draw (s5n16) -- (s5n17);
  \draw (s5n11) -- (s5n16);
  \draw (s5n11) -- (s5n18);
  \draw (s5n10) -- (s5n11);
  \draw (s5n20) -- (s5n21);
  \draw (s5n19) -- (s5n20);
  \draw (s5n19) -- (s5n22);
  \draw (s5n10) -- (s5n19);
  \draw (s5n23) -- (s5n24);
  \draw (s5n10) -- (s5n23);
  \draw (s5n10) -- (s5n25);
  \draw[->] (s5n0) -- (s5n8);
  \draw[->] (s5n8) -- (s5n10);
  \node[head, anchor=east] (s5h0) at (1.094,-16.250) {$\mathit{head}[H]$};
  \draw[->] (s5h0) -- (s5n0);
  \node[ptr, anchor=south] at (s5n0.north) {\textit{prev-x}};
  \node[ptr, anchor=south] at (s5n8.north) {\textit{x}};
  \node[ptr, anchor=south] at (s5n10.north) {\textit{next-x}};
\end{tikzpicture}
\end{widget}
\caption{L'unione di due heap binomiali con $7$ e $19$ nodi. (a) I due heap $H_1$ e $H_2$. (b) La lista prodotta da \BinomialHeapMerge{}, con i puntatori all'inizio del ciclo. (c)--(f) La lista all'inizio delle iterazioni successive, dopo i casi 3, 2, 4 e 3, con in evidenza le radici appena collegate. Nell'ultima iterazione si applica il caso 1 e il ciclo termina.}
\label{fig:hb-unione}
\end{figure}

\paragraph{Costo.}
Siano $n_1$ e $n_2$ i nodi di $H_1$ e $H_2$, e $n = n_1 + n_2$. I due heap hanno al più $\lfloor \lg n_1 \rfloor + 1$ e $\lfloor \lg n_2 \rfloor + 1$ alberi, quindi la lista prodotta da \BinomialHeapMerge{} contiene al più
\[
\lfloor \lg n_1 \rfloor + \lfloor \lg n_2 \rfloor + 2 \le 2 \lfloor \lg n \rfloor + 2
\]
radici, dove la disuguaglianza vale in quanto $n_1 \le n$ e $n_2 \le n$. La fusione, quindi, costa $O(\lg n)$.

Per il ciclo \textbf{while} osserviamo che a ogni iterazione il puntatore $x$ avanza di una posizione (casi 1 e 2), oppure una radice viene tolta dalla lista (casi 3 e 4). Entrambe le cose possono accadere al più tante volte quante sono le radici, e ogni iterazione richiede tempo costante. Il costo dell'unione è quindi:
\[
O(\lg n).
\]

\subsection{Inserimento}\label{sec:hb-insert}
Per inserire un nodo $x$ si costruisce un heap $H'$ che contiene soltanto $x$, ovvero un $B_0$, e lo si unisce a $H$ (algoritmo~\ref{alg:hb-insert}). Il costo è quello dell'unione, $O(\lg n)$, dove $n$ è il numero di nodi di $H$.

\begin{algorithm}[htbp]
\caption{L'inserimento del nodo $x$, come unione con un heap che contiene soltanto $x$.}\label{alg:hb-insert}
\begin{algorithmic}[1]
\Procedure{Binomial-Heap-Insert}{$H,x$}
  \State $H' \gets$ \Call{Make-Binomial-Heap}{}
  \State $p[x] \gets \Nil$
  \State $\mathit{child}[x] \gets \Nil$
  \State $\mathit{sibling}[x] \gets \Nil$
  \State $\mathit{degree}[x] \gets 0$
  \State $\mathit{head}[H'] \gets x$
  \State $H \gets$ \Call{Binomial-Heap-Union}{$H,H'$}
\EndProcedure
\end{algorithmic}
\end{algorithm}

\subsection{Estrazione del minimo}\label{sec:hb-extract}
\BinomialHeapExtractMin$(H)$ toglie dall'heap il nodo con la chiave minima, e ne restituisce il puntatore (algoritmo~\ref{alg:hb-extract}).

\begin{algorithm}[htbp]
\caption{L'estrazione del nodo con la chiave minima: i figli della radice tolta formano un nuovo heap, che si unisce a quello di partenza.}\label{alg:hb-extract}
\begin{algorithmic}[1]
\Procedure{Binomial-Heap-Extract-Min}{$H$}
  \State trova la radice $x$ con la chiave minima nella lista delle radici di $H$
  \State togli $x$ dalla lista delle radici di $H$
  \State $H' \gets$ \Call{Make-Binomial-Heap}{}
  \State inverti l'ordine della lista dei figli di $x$
  \State $\mathit{head}[H'] \gets$ la prima radice della lista invertita
  \State $H \gets$ \Call{Binomial-Heap-Union}{$H,H'$}
  \State \Return $x$
\EndProcedure
\end{algorithmic}
\end{algorithm}

Le prime due righe trovano e tolgono la radice $x$ con la chiave minima, che per la proprietà di min-heap è il minimo dell'intero heap. I figli di $x$, però, vanno rimessi nell'heap. Se $x$ è la radice di un $B_k$, per la proprietà 4(c) i suoi figli sono, da sinistra, le radici di $B_{k-1}, \ldots, B_1, B_0$, ovvero sono in ordine di grado decrescente. Invertendo la loro lista si ottiene quindi la lista delle radici, ordinata per grado crescente, di un nuovo heap binomiale $H'$, che si unisce a $H$.

\paragraph{Costo.}
La ricerca di $x$ costa $O(\lg n)$, come \BinomialHeapMinimum{}, e la creazione di $H'$ costa $O(1)$. Per il corollario~\ref{cor:hb-grado-heap} la radice $x$ ha al più $\lfloor \lg n \rfloor$ figli, quindi anche l'inversione della lista dei figli costa $O(\lg n)$. Infine l'unione costa $O(\lg n)$, e il costo complessivo è:
\[
O(\lg n).
\]

\paragraph{Esempio.}
Consideriamo l'heap di figura~\ref{fig:hb-estrazione}a, con $n = 22 = (10110)_2$ nodi, e quindi con gli alberi $B_1$, $B_2$ e $B_4$. Il minimo è la radice con chiave 1, la radice del $B_4$. Togliendola, $H$ resta con $B_1$ e $B_2$, ovvero con $6 = (00110)_2$ nodi. I figli di 1, in ordine inverso, formano $H'$, con $B_0$, $B_1$, $B_2$ e $B_3$, ovvero con $15 = (01111)_2$ nodi (figura~\ref{fig:hb-estrazione}b).

L'unione dei due heap produce l'heap di figura~\ref{fig:hb-estrazione}c, con
\[
(00110)_2 + (01111)_2 = (10101)_2 = 21
\]
nodi, ovvero con $B_0$, $B_2$ e $B_4$. Durante l'unione si applicano, nell'ordine, i casi 1, 4, 2, 3 e 4.

\begin{figure}[htbp]
\centering
\begin{widget}{binomial-heap}[heap={(37 (41)) (10 (28 (77)) (13)) (1 (6 (8 (11 (27)) (17)) (14 (38)) (29)) (16 (26 (42)) (23)) (12 (18)) (25))}, ops={extract-min}, actions={insert, extract-min, decrease-key, delete}]
\begin{tikzpicture}[binheap, >=Stealth]
  % (a) 
  \node[lab, anchor=west] at (-1.05,0.000) {(a)};
  \node[hn] (s0n1) at (1.530,-0.600) {41};
  \node[hn] (s0n0) at (1.530,0.000) {37};
  \node[hn] (s0n4) at (2.510,-1.200) {77};
  \node[hn] (s0n3) at (2.510,-0.600) {28};
  \node[hn] (s0n5) at (3.070,-0.600) {13};
  \node[hn] (s0n2) at (3.070,0.000) {10};
  \node[hn] (s0n10) at (4.050,-2.400) {27};
  \node[hn] (s0n9) at (4.050,-1.800) {11};
  \node[hn] (s0n11) at (4.610,-1.800) {17};
  \node[hn] (s0n8) at (4.610,-1.200) {8};
  \node[hn] (s0n13) at (5.170,-1.800) {38};
  \node[hn] (s0n12) at (5.170,-1.200) {14};
  \node[hn] (s0n14) at (5.730,-1.200) {29};
  \node[hn] (s0n7) at (5.730,-0.600) {6};
  \node[hn] (s0n17) at (6.290,-1.800) {42};
  \node[hn] (s0n16) at (6.290,-1.200) {26};
  \node[hn] (s0n18) at (6.850,-1.200) {23};
  \node[hn] (s0n15) at (6.850,-0.600) {16};
  \node[hn] (s0n20) at (7.410,-1.200) {18};
  \node[hn] (s0n19) at (7.410,-0.600) {12};
  \node[hn] (s0n21) at (7.970,-0.600) {25};
  \node[hl] (s0n6) at (7.970,0.000) {1};
  \draw (s0n0) -- (s0n1);
  \draw (s0n3) -- (s0n4);
  \draw (s0n2) -- (s0n3);
  \draw (s0n2) -- (s0n5);
  \draw (s0n9) -- (s0n10);
  \draw (s0n8) -- (s0n9);
  \draw (s0n8) -- (s0n11);
  \draw (s0n7) -- (s0n8);
  \draw (s0n12) -- (s0n13);
  \draw (s0n7) -- (s0n12);
  \draw (s0n7) -- (s0n14);
  \draw (s0n6) -- (s0n7);
  \draw (s0n16) -- (s0n17);
  \draw (s0n15) -- (s0n16);
  \draw (s0n15) -- (s0n18);
  \draw (s0n6) -- (s0n15);
  \draw (s0n19) -- (s0n20);
  \draw (s0n6) -- (s0n19);
  \draw (s0n6) -- (s0n21);
  \draw[->] (s0n0) -- (s0n2);
  \draw[->] (s0n2) -- (s0n6);
  \node[head, anchor=east] (s0h0) at (1.094,0.000) {$\mathit{head}[H]$};
  \draw[->] (s0h0) -- (s0n0);
  \node[ptr, anchor=south] at (s0n6.north) {\textit{x}};
  % (b) 
  \node[lab, anchor=west] at (-1.05,-3.150) {(b)};
  \node[hn] (s1n1) at (1.530,-3.750) {41};
  \node[hn] (s1n0) at (1.530,-3.150) {37};
  \node[hn] (s1n4) at (2.510,-4.350) {77};
  \node[hn] (s1n3) at (2.510,-3.750) {28};
  \node[hn] (s1n5) at (3.070,-3.750) {13};
  \node[hn] (s1n2) at (3.070,-3.150) {10};
  \node[hn] (s1n21) at (5.248,-3.150) {25};
  \node[hn] (s1n20) at (6.228,-3.750) {18};
  \node[hn] (s1n19) at (6.228,-3.150) {12};
  \node[hn] (s1n17) at (7.208,-4.350) {42};
  \node[hn] (s1n16) at (7.208,-3.750) {26};
  \node[hn] (s1n18) at (7.768,-3.750) {23};
  \node[hn] (s1n15) at (7.768,-3.150) {16};
  \node[hn] (s1n10) at (8.748,-4.950) {27};
  \node[hn] (s1n9) at (8.748,-4.350) {11};
  \node[hn] (s1n11) at (9.308,-4.350) {17};
  \node[hn] (s1n8) at (9.308,-3.750) {8};
  \node[hn] (s1n13) at (9.868,-4.350) {38};
  \node[hn] (s1n12) at (9.868,-3.750) {14};
  \node[hn] (s1n14) at (10.428,-3.750) {29};
  \node[hn] (s1n7) at (10.428,-3.150) {6};
  \draw (s1n0) -- (s1n1);
  \draw (s1n3) -- (s1n4);
  \draw (s1n2) -- (s1n3);
  \draw (s1n2) -- (s1n5);
  \draw (s1n19) -- (s1n20);
  \draw (s1n16) -- (s1n17);
  \draw (s1n15) -- (s1n16);
  \draw (s1n15) -- (s1n18);
  \draw (s1n9) -- (s1n10);
  \draw (s1n8) -- (s1n9);
  \draw (s1n8) -- (s1n11);
  \draw (s1n7) -- (s1n8);
  \draw (s1n12) -- (s1n13);
  \draw (s1n7) -- (s1n12);
  \draw (s1n7) -- (s1n14);
  \draw[->] (s1n0) -- (s1n2);
  \draw[->] (s1n21) -- (s1n19);
  \draw[->] (s1n19) -- (s1n15);
  \draw[->] (s1n15) -- (s1n7);
  \node[head, anchor=east] (s1h0) at (1.094,-3.150) {$\mathit{head}[H]$};
  \draw[->] (s1h0) -- (s1n0);
  \node[head, anchor=east] (s1h1) at (4.812,-3.150) {$\mathit{head}[H']$};
  \draw[->] (s1h1) -- (s1n21);
  % (c) 
  \node[lab, anchor=west] at (-1.05,-5.700) {(c)};
  \node[hn] (s2n21) at (1.530,-5.700) {25};
  \node[hn] (s2n1) at (2.510,-6.900) {41};
  \node[hn] (s2n0) at (2.510,-6.300) {37};
  \node[hn] (s2n20) at (3.070,-6.300) {18};
  \node[hn] (s2n19) at (3.070,-5.700) {12};
  \node[hn] (s2n17) at (4.050,-8.100) {42};
  \node[hn] (s2n16) at (4.050,-7.500) {26};
  \node[hn] (s2n18) at (4.610,-7.500) {23};
  \node[hn] (s2n15) at (4.610,-6.900) {16};
  \node[hn] (s2n4) at (5.170,-7.500) {77};
  \node[hn] (s2n3) at (5.170,-6.900) {28};
  \node[hn] (s2n5) at (5.730,-6.900) {13};
  \node[hn] (s2n2) at (5.730,-6.300) {10};
  \node[hn] (s2n10) at (6.290,-7.500) {27};
  \node[hn] (s2n9) at (6.290,-6.900) {11};
  \node[hn] (s2n11) at (6.850,-6.900) {17};
  \node[hn] (s2n8) at (6.850,-6.300) {8};
  \node[hn] (s2n13) at (7.410,-6.900) {38};
  \node[hn] (s2n12) at (7.410,-6.300) {14};
  \node[hn] (s2n14) at (7.970,-6.300) {29};
  \node[hn] (s2n7) at (7.970,-5.700) {6};
  \draw (s2n0) -- (s2n1);
  \draw (s2n19) -- (s2n0);
  \draw (s2n19) -- (s2n20);
  \draw (s2n16) -- (s2n17);
  \draw (s2n15) -- (s2n16);
  \draw (s2n15) -- (s2n18);
  \draw (s2n2) -- (s2n15);
  \draw (s2n3) -- (s2n4);
  \draw (s2n2) -- (s2n3);
  \draw (s2n2) -- (s2n5);
  \draw (s2n7) -- (s2n2);
  \draw (s2n9) -- (s2n10);
  \draw (s2n8) -- (s2n9);
  \draw (s2n8) -- (s2n11);
  \draw (s2n7) -- (s2n8);
  \draw (s2n12) -- (s2n13);
  \draw (s2n7) -- (s2n12);
  \draw (s2n7) -- (s2n14);
  \draw[->] (s2n21) -- (s2n19);
  \draw[->] (s2n19) -- (s2n7);
  \node[head, anchor=east] (s2h0) at (1.094,-5.700) {$\mathit{head}[H]$};
  \draw[->] (s2h0) -- (s2n21);
\end{tikzpicture}
\end{widget}
\caption{L'estrazione del minimo da un heap binomiale con $22$ nodi. (a) Il minimo è la radice $x$ con chiave 1, la radice di un $B_4$. (b) Tolto $x$, i suoi figli in ordine inverso formano la lista delle radici di $H'$. (c) L'heap ottenuto dall'unione di $H$ e $H'$, con $21$ nodi.}
\label{fig:hb-estrazione}
\end{figure}

\subsection{Decremento di una chiave}\label{sec:hb-decrease}
\BinomialHeapDecreaseKey$(H,x,k)$ assegna al nodo $x$ la chiave $k$, che non deve essere maggiore di quella attuale (algoritmo~\ref{alg:hb-decrease}).

\begin{algorithm}[htbp]
\caption{Il decremento della chiave del nodo $x$: la nuova chiave risale verso la radice finché è minore di quella del padre.}\label{alg:hb-decrease}
\begin{algorithmic}[1]
\Procedure{Binomial-Heap-Decrease-Key}{$H,x,k$}
  \If{$k > \mathit{key}[x]$}
    \State \textbf{error} ``la nuova chiave è maggiore di quella attuale''
  \EndIf
  \State $\mathit{key}[x] \gets k$
  \State $y \gets x$
  \State $z \gets p[y]$
  \While{$z \ne \Nil$ \textbf{and} $\mathit{key}[y] < \mathit{key}[z]$}
    \State scambia $\mathit{key}[y]$ e $\mathit{key}[z]$ \Comment{insieme alle informazioni satellite}
    \State $y \gets z$
    \State $z \gets p[y]$
  \EndWhile
\EndProcedure
\end{algorithmic}
\end{algorithm}

Diminuendo la chiave di $x$, la proprietà di min-heap può non valere più tra $x$ e suo padre. La si ripristina facendo risalire la chiave. Finché la chiave di $y$ è minore di quella del padre $z$, si scambiano le due chiavi, insieme alle informazioni satellite, e si sale di un livello. Il ciclo termina quando $y$ è una radice, oppure quando la sua chiave non è più minore di quella del padre.

\paragraph{Costo.}
Ogni iterazione costa $O(1)$ e fa salire $y$ di un livello, quindi le iterazioni sono al più tante quanto la profondità di $x$. Il nodo $x$ appartiene a un albero $B_k$ con $2^k \le n$, che per la proprietà 2 ha altezza $k \le \lfloor \lg n \rfloor$. Il costo è quindi $O(\lg n)$.

\paragraph{Esempio.}
Nella figura~\ref{fig:hb-decremento} la chiave 26 dell'heap ottenuto nell'esempio precedente diventa 7. La chiave risale due volte, scambiandosi prima con 16 e poi con 10, e si ferma sotto la radice, in quanto $7 \ge 6$.

\begin{figure}[htbp]
\centering
\begin{widget}{binomial-heap}[heap={(25) (12 (37 (41)) (18)) (6 (10 (16 (26 (42)) (23)) (28 (77)) (13)) (8 (11 (27)) (17)) (14 (38)) (29))}, ops={decrease-key 26 7}, actions={decrease-key, delete}]
\begin{tikzpicture}[binheap, >=Stealth]
  % (a) 
  \node[lab, anchor=west] at (-1.05,0.000) {(a)};
  \node[hn] (s0n21) at (1.530,0.000) {25};
  \node[hn] (s0n1) at (2.510,-1.200) {41};
  \node[hn] (s0n0) at (2.510,-0.600) {37};
  \node[hn] (s0n20) at (3.070,-0.600) {18};
  \node[hn] (s0n19) at (3.070,0.000) {12};
  \node[hn] (s0n17) at (4.050,-2.400) {42};
  \node[hl] (s0n16) at (4.050,-1.800) {7};
  \node[hn] (s0n18) at (4.610,-1.800) {23};
  \node[hn] (s0n15) at (4.610,-1.200) {16};
  \node[hn] (s0n4) at (5.170,-1.800) {77};
  \node[hn] (s0n3) at (5.170,-1.200) {28};
  \node[hn] (s0n5) at (5.730,-1.200) {13};
  \node[hn] (s0n2) at (5.730,-0.600) {10};
  \node[hn] (s0n10) at (6.290,-1.800) {27};
  \node[hn] (s0n9) at (6.290,-1.200) {11};
  \node[hn] (s0n11) at (6.850,-1.200) {17};
  \node[hn] (s0n8) at (6.850,-0.600) {8};
  \node[hn] (s0n13) at (7.410,-1.200) {38};
  \node[hn] (s0n12) at (7.410,-0.600) {14};
  \node[hn] (s0n14) at (7.970,-0.600) {29};
  \node[hn] (s0n7) at (7.970,0.000) {6};
  \draw (s0n0) -- (s0n1);
  \draw (s0n19) -- (s0n0);
  \draw (s0n19) -- (s0n20);
  \draw (s0n16) -- (s0n17);
  \draw (s0n15) -- (s0n16);
  \draw (s0n15) -- (s0n18);
  \draw (s0n2) -- (s0n15);
  \draw (s0n3) -- (s0n4);
  \draw (s0n2) -- (s0n3);
  \draw (s0n2) -- (s0n5);
  \draw (s0n7) -- (s0n2);
  \draw (s0n9) -- (s0n10);
  \draw (s0n8) -- (s0n9);
  \draw (s0n8) -- (s0n11);
  \draw (s0n7) -- (s0n8);
  \draw (s0n12) -- (s0n13);
  \draw (s0n7) -- (s0n12);
  \draw (s0n7) -- (s0n14);
  \draw[->] (s0n21) -- (s0n19);
  \draw[->] (s0n19) -- (s0n7);
  \node[head, anchor=east] (s0h0) at (1.094,0.000) {$\mathit{head}[H]$};
  \draw[->] (s0h0) -- (s0n21);
  \node[ptr, anchor=east] at (s0n16.west) {\textit{y}};
  \node[ptr, anchor=east] at (s0n15.west) {\textit{z}};
  % (b) 
  \node[lab, anchor=west] at (-1.05,-3.450) {(b)};
  \node[hn] (s1n21) at (1.530,-3.450) {25};
  \node[hn] (s1n1) at (2.510,-4.650) {41};
  \node[hn] (s1n0) at (2.510,-4.050) {37};
  \node[hn] (s1n20) at (3.070,-4.050) {18};
  \node[hn] (s1n19) at (3.070,-3.450) {12};
  \node[hn] (s1n17) at (4.050,-5.850) {42};
  \node[hn] (s1n16) at (4.050,-5.250) {16};
  \node[hn] (s1n18) at (4.610,-5.250) {23};
  \node[hl] (s1n15) at (4.610,-4.650) {7};
  \node[hn] (s1n4) at (5.170,-5.250) {77};
  \node[hn] (s1n3) at (5.170,-4.650) {28};
  \node[hn] (s1n5) at (5.730,-4.650) {13};
  \node[hn] (s1n2) at (5.730,-4.050) {10};
  \node[hn] (s1n10) at (6.290,-5.250) {27};
  \node[hn] (s1n9) at (6.290,-4.650) {11};
  \node[hn] (s1n11) at (6.850,-4.650) {17};
  \node[hn] (s1n8) at (6.850,-4.050) {8};
  \node[hn] (s1n13) at (7.410,-4.650) {38};
  \node[hn] (s1n12) at (7.410,-4.050) {14};
  \node[hn] (s1n14) at (7.970,-4.050) {29};
  \node[hn] (s1n7) at (7.970,-3.450) {6};
  \draw (s1n0) -- (s1n1);
  \draw (s1n19) -- (s1n0);
  \draw (s1n19) -- (s1n20);
  \draw (s1n16) -- (s1n17);
  \draw (s1n15) -- (s1n16);
  \draw (s1n15) -- (s1n18);
  \draw (s1n2) -- (s1n15);
  \draw (s1n3) -- (s1n4);
  \draw (s1n2) -- (s1n3);
  \draw (s1n2) -- (s1n5);
  \draw (s1n7) -- (s1n2);
  \draw (s1n9) -- (s1n10);
  \draw (s1n8) -- (s1n9);
  \draw (s1n8) -- (s1n11);
  \draw (s1n7) -- (s1n8);
  \draw (s1n12) -- (s1n13);
  \draw (s1n7) -- (s1n12);
  \draw (s1n7) -- (s1n14);
  \draw[->] (s1n21) -- (s1n19);
  \draw[->] (s1n19) -- (s1n7);
  \node[head, anchor=east] (s1h0) at (1.094,-3.450) {$\mathit{head}[H]$};
  \draw[->] (s1h0) -- (s1n21);
  \node[ptr, anchor=east] at (s1n15.west) {\textit{y}};
  \node[ptr, anchor=east] at (s1n2.west) {\textit{z}};
  % (c) 
  \node[lab, anchor=west] at (-1.05,-6.900) {(c)};
  \node[hn] (s2n21) at (1.530,-6.900) {25};
  \node[hn] (s2n1) at (2.510,-8.100) {41};
  \node[hn] (s2n0) at (2.510,-7.500) {37};
  \node[hn] (s2n20) at (3.070,-7.500) {18};
  \node[hn] (s2n19) at (3.070,-6.900) {12};
  \node[hn] (s2n17) at (4.050,-9.300) {42};
  \node[hn] (s2n16) at (4.050,-8.700) {16};
  \node[hn] (s2n18) at (4.610,-8.700) {23};
  \node[hn] (s2n15) at (4.610,-8.100) {10};
  \node[hn] (s2n4) at (5.170,-8.700) {77};
  \node[hn] (s2n3) at (5.170,-8.100) {28};
  \node[hn] (s2n5) at (5.730,-8.100) {13};
  \node[hl] (s2n2) at (5.730,-7.500) {7};
  \node[hn] (s2n10) at (6.290,-8.700) {27};
  \node[hn] (s2n9) at (6.290,-8.100) {11};
  \node[hn] (s2n11) at (6.850,-8.100) {17};
  \node[hn] (s2n8) at (6.850,-7.500) {8};
  \node[hn] (s2n13) at (7.410,-8.100) {38};
  \node[hn] (s2n12) at (7.410,-7.500) {14};
  \node[hn] (s2n14) at (7.970,-7.500) {29};
  \node[hn] (s2n7) at (7.970,-6.900) {6};
  \draw (s2n0) -- (s2n1);
  \draw (s2n19) -- (s2n0);
  \draw (s2n19) -- (s2n20);
  \draw (s2n16) -- (s2n17);
  \draw (s2n15) -- (s2n16);
  \draw (s2n15) -- (s2n18);
  \draw (s2n2) -- (s2n15);
  \draw (s2n3) -- (s2n4);
  \draw (s2n2) -- (s2n3);
  \draw (s2n2) -- (s2n5);
  \draw (s2n7) -- (s2n2);
  \draw (s2n9) -- (s2n10);
  \draw (s2n8) -- (s2n9);
  \draw (s2n8) -- (s2n11);
  \draw (s2n7) -- (s2n8);
  \draw (s2n12) -- (s2n13);
  \draw (s2n7) -- (s2n12);
  \draw (s2n7) -- (s2n14);
  \draw[->] (s2n21) -- (s2n19);
  \draw[->] (s2n19) -- (s2n7);
  \node[head, anchor=east] (s2h0) at (1.094,-6.900) {$\mathit{head}[H]$};
  \draw[->] (s2h0) -- (s2n21);
  \node[ptr, anchor=east] at (s2n2.west) {\textit{y}};
  \node[ptr, anchor=south] at (s2n7.north) {\textit{z}};
\end{tikzpicture}
\end{widget}
\caption{Il decremento della chiave 26 a 7. (a) La nuova chiave è minore di quella del padre $z$. (b) Dopo il primo scambio la chiave 7 è salita di un livello, ed è ancora minore di quella del padre. (c) Dopo il secondo scambio $\mathit{key}[y] \ge \mathit{key}[z]$, e il ciclo termina.}
\label{fig:hb-decremento}
\end{figure}

\subsection{Cancellazione}\label{sec:hb-delete}
Per cancellare un nodo $x$ si riutilizzano le due operazioni precedenti (algoritmo~\ref{alg:hb-delete}): prima si porta la chiave di $x$ a $-\infty$, poi si estrae il minimo.

\begin{algorithm}[htbp]
\caption{La cancellazione del nodo $x$: con la chiave $-\infty$, $x$ diventa il minimo dell'heap.}\label{alg:hb-delete}
\begin{algorithmic}[1]
\Procedure{Binomial-Heap-Delete}{$H,x$}
  \State \Call{Binomial-Heap-Decrease-Key}{$H,x,-\infty$}
  \State \Call{Binomial-Heap-Extract-Min}{$H$}
\EndProcedure
\end{algorithmic}
\end{algorithm}

La chiave $-\infty$ è minore di tutte le altre, quindi risale fino alla radice del suo albero e diventa il minimo dell'heap. \BinomialHeapExtractMin{} rimuove quindi proprio il nodo che la contiene. Entrambe le operazioni costano $O(\lg n)$, e quindi anche la cancellazione.

\paragraph{Esempio.}
Nella figura~\ref{fig:hb-cancellazione} si cancella il nodo con chiave 77 dall'heap ottenuto nell'esempio precedente. La chiave $-\infty$ risale fino alla radice del $B_4$ (b), che viene estratta. I suoi figli in ordine inverso formano $H'$ (c), e l'unione produce un heap con $20 = (10100)_2$ nodi, ovvero con $B_2$ e $B_4$ (d).

\begin{figure}[htbp]
\centering
\begin{widget}{binomial-heap}[heap={(25) (12 (37 (41)) (18)) (6 (7 (10 (16 (42)) (23)) (28 (77)) (13)) (8 (11 (27)) (17)) (14 (38)) (29))}, ops={delete 77}, actions={delete, decrease-key}]
\begin{tikzpicture}[binheap, >=Stealth]
  % (a) 
  \node[lab, anchor=west] at (-1.05,0.000) {(a)};
  \node[hn] (s0n21) at (1.530,0.000) {25};
  \node[hn] (s0n1) at (2.510,-1.200) {41};
  \node[hn] (s0n0) at (2.510,-0.600) {37};
  \node[hn] (s0n20) at (3.070,-0.600) {18};
  \node[hn] (s0n19) at (3.070,0.000) {12};
  \node[hn] (s0n17) at (4.050,-2.400) {42};
  \node[hn] (s0n16) at (4.050,-1.800) {16};
  \node[hn] (s0n18) at (4.610,-1.800) {23};
  \node[hn] (s0n15) at (4.610,-1.200) {10};
  \node[hl] (s0n4) at (5.170,-1.800) {77};
  \node[hn] (s0n3) at (5.170,-1.200) {28};
  \node[hn] (s0n5) at (5.730,-1.200) {13};
  \node[hn] (s0n2) at (5.730,-0.600) {7};
  \node[hn] (s0n10) at (6.290,-1.800) {27};
  \node[hn] (s0n9) at (6.290,-1.200) {11};
  \node[hn] (s0n11) at (6.850,-1.200) {17};
  \node[hn] (s0n8) at (6.850,-0.600) {8};
  \node[hn] (s0n13) at (7.410,-1.200) {38};
  \node[hn] (s0n12) at (7.410,-0.600) {14};
  \node[hn] (s0n14) at (7.970,-0.600) {29};
  \node[hn] (s0n7) at (7.970,0.000) {6};
  \draw (s0n0) -- (s0n1);
  \draw (s0n19) -- (s0n0);
  \draw (s0n19) -- (s0n20);
  \draw (s0n16) -- (s0n17);
  \draw (s0n15) -- (s0n16);
  \draw (s0n15) -- (s0n18);
  \draw (s0n2) -- (s0n15);
  \draw (s0n3) -- (s0n4);
  \draw (s0n2) -- (s0n3);
  \draw (s0n2) -- (s0n5);
  \draw (s0n7) -- (s0n2);
  \draw (s0n9) -- (s0n10);
  \draw (s0n8) -- (s0n9);
  \draw (s0n8) -- (s0n11);
  \draw (s0n7) -- (s0n8);
  \draw (s0n12) -- (s0n13);
  \draw (s0n7) -- (s0n12);
  \draw (s0n7) -- (s0n14);
  \draw[->] (s0n21) -- (s0n19);
  \draw[->] (s0n19) -- (s0n7);
  \node[head, anchor=east] (s0h0) at (1.094,0.000) {$\mathit{head}[H]$};
  \draw[->] (s0h0) -- (s0n21);
  % (b) 
  \node[lab, anchor=west] at (-1.05,-3.150) {(b)};
  \node[hn] (s1n21) at (1.530,-3.150) {25};
  \node[hn] (s1n1) at (2.510,-4.350) {41};
  \node[hn] (s1n0) at (2.510,-3.750) {37};
  \node[hn] (s1n20) at (3.070,-3.750) {18};
  \node[hn] (s1n19) at (3.070,-3.150) {12};
  \node[hn] (s1n17) at (4.050,-5.550) {42};
  \node[hn] (s1n16) at (4.050,-4.950) {16};
  \node[hn] (s1n18) at (4.610,-4.950) {23};
  \node[hn] (s1n15) at (4.610,-4.350) {10};
  \node[hn] (s1n4) at (5.170,-4.950) {28};
  \node[hn] (s1n3) at (5.170,-4.350) {7};
  \node[hn] (s1n5) at (5.730,-4.350) {13};
  \node[hn] (s1n2) at (5.730,-3.750) {6};
  \node[hn] (s1n10) at (6.290,-4.950) {27};
  \node[hn] (s1n9) at (6.290,-4.350) {11};
  \node[hn] (s1n11) at (6.850,-4.350) {17};
  \node[hn] (s1n8) at (6.850,-3.750) {8};
  \node[hn] (s1n13) at (7.410,-4.350) {38};
  \node[hn] (s1n12) at (7.410,-3.750) {14};
  \node[hn] (s1n14) at (7.970,-3.750) {29};
  \node[hl] (s1n7) at (7.970,-3.150) {$-\infty$};
  \draw (s1n0) -- (s1n1);
  \draw (s1n19) -- (s1n0);
  \draw (s1n19) -- (s1n20);
  \draw (s1n16) -- (s1n17);
  \draw (s1n15) -- (s1n16);
  \draw (s1n15) -- (s1n18);
  \draw (s1n2) -- (s1n15);
  \draw (s1n3) -- (s1n4);
  \draw (s1n2) -- (s1n3);
  \draw (s1n2) -- (s1n5);
  \draw (s1n7) -- (s1n2);
  \draw (s1n9) -- (s1n10);
  \draw (s1n8) -- (s1n9);
  \draw (s1n8) -- (s1n11);
  \draw (s1n7) -- (s1n8);
  \draw (s1n12) -- (s1n13);
  \draw (s1n7) -- (s1n12);
  \draw (s1n7) -- (s1n14);
  \draw[->] (s1n21) -- (s1n19);
  \draw[->] (s1n19) -- (s1n7);
  \node[head, anchor=east] (s1h0) at (1.094,-3.150) {$\mathit{head}[H]$};
  \draw[->] (s1h0) -- (s1n21);
  % (c) 
  \node[lab, anchor=west] at (-1.05,-6.300) {(c)};
  \node[hn] (s2n21) at (1.530,-6.300) {25};
  \node[hn] (s2n1) at (2.510,-7.500) {41};
  \node[hn] (s2n0) at (2.510,-6.900) {37};
  \node[hn] (s2n20) at (3.070,-6.900) {18};
  \node[hn] (s2n19) at (3.070,-6.300) {12};
  \node[hn] (s2n14) at (5.248,-6.300) {29};
  \node[hn] (s2n13) at (6.228,-6.900) {38};
  \node[hn] (s2n12) at (6.228,-6.300) {14};
  \node[hn] (s2n10) at (7.208,-7.500) {27};
  \node[hn] (s2n9) at (7.208,-6.900) {11};
  \node[hn] (s2n11) at (7.768,-6.900) {17};
  \node[hn] (s2n8) at (7.768,-6.300) {8};
  \node[hn] (s2n17) at (8.748,-8.100) {42};
  \node[hn] (s2n16) at (8.748,-7.500) {16};
  \node[hn] (s2n18) at (9.308,-7.500) {23};
  \node[hn] (s2n15) at (9.308,-6.900) {10};
  \node[hn] (s2n4) at (9.868,-7.500) {28};
  \node[hn] (s2n3) at (9.868,-6.900) {7};
  \node[hn] (s2n5) at (10.428,-6.900) {13};
  \node[hn] (s2n2) at (10.428,-6.300) {6};
  \draw (s2n0) -- (s2n1);
  \draw (s2n19) -- (s2n0);
  \draw (s2n19) -- (s2n20);
  \draw (s2n12) -- (s2n13);
  \draw (s2n9) -- (s2n10);
  \draw (s2n8) -- (s2n9);
  \draw (s2n8) -- (s2n11);
  \draw (s2n16) -- (s2n17);
  \draw (s2n15) -- (s2n16);
  \draw (s2n15) -- (s2n18);
  \draw (s2n2) -- (s2n15);
  \draw (s2n3) -- (s2n4);
  \draw (s2n2) -- (s2n3);
  \draw (s2n2) -- (s2n5);
  \draw[->] (s2n21) -- (s2n19);
  \draw[->] (s2n14) -- (s2n12);
  \draw[->] (s2n12) -- (s2n8);
  \draw[->] (s2n8) -- (s2n2);
  \node[head, anchor=east] (s2h0) at (1.094,-6.300) {$\mathit{head}[H]$};
  \draw[->] (s2h0) -- (s2n21);
  \node[head, anchor=east] (s2h1) at (4.812,-6.300) {$\mathit{head}[H']$};
  \draw[->] (s2h1) -- (s2n14);
  % (d) 
  \node[lab, anchor=west] at (-1.05,-8.850) {(d)};
  \node[hn] (s3n14) at (1.530,-10.050) {29};
  \node[hn] (s3n21) at (1.530,-9.450) {25};
  \node[hn] (s3n13) at (2.090,-9.450) {38};
  \node[hn] (s3n12) at (2.090,-8.850) {14};
  \node[hn] (s3n1) at (3.070,-11.250) {41};
  \node[hn] (s3n0) at (3.070,-10.650) {37};
  \node[hn] (s3n20) at (3.630,-10.650) {18};
  \node[hn] (s3n19) at (3.630,-10.050) {12};
  \node[hn] (s3n10) at (4.190,-10.650) {27};
  \node[hn] (s3n9) at (4.190,-10.050) {11};
  \node[hn] (s3n11) at (4.750,-10.050) {17};
  \node[hn] (s3n8) at (4.750,-9.450) {8};
  \node[hn] (s3n17) at (5.310,-10.650) {42};
  \node[hn] (s3n16) at (5.310,-10.050) {16};
  \node[hn] (s3n18) at (5.870,-10.050) {23};
  \node[hn] (s3n15) at (5.870,-9.450) {10};
  \node[hn] (s3n4) at (6.430,-10.050) {28};
  \node[hn] (s3n3) at (6.430,-9.450) {7};
  \node[hn] (s3n5) at (6.990,-9.450) {13};
  \node[hn] (s3n2) at (6.990,-8.850) {6};
  \draw (s3n21) -- (s3n14);
  \draw (s3n12) -- (s3n21);
  \draw (s3n12) -- (s3n13);
  \draw (s3n0) -- (s3n1);
  \draw (s3n19) -- (s3n0);
  \draw (s3n19) -- (s3n20);
  \draw (s3n8) -- (s3n19);
  \draw (s3n9) -- (s3n10);
  \draw (s3n8) -- (s3n9);
  \draw (s3n8) -- (s3n11);
  \draw (s3n2) -- (s3n8);
  \draw (s3n16) -- (s3n17);
  \draw (s3n15) -- (s3n16);
  \draw (s3n15) -- (s3n18);
  \draw (s3n2) -- (s3n15);
  \draw (s3n3) -- (s3n4);
  \draw (s3n2) -- (s3n3);
  \draw (s3n2) -- (s3n5);
  \draw[->] (s3n12) -- (s3n2);
  \node[head, anchor=east] (s3h0) at (1.094,-8.850) {$\mathit{head}[H]$};
  \draw[->] (s3h0) -- (s3n12);
\end{tikzpicture}
\end{widget}
\caption{La cancellazione del nodo con chiave 77. (a) L'heap di partenza. (b) Dopo il decremento a $-\infty$, la chiave è salita fino alla radice del $B_4$. (c) Estratto il minimo, i figli della radice tolta formano $H'$. (d) L'heap ottenuto dall'unione, con $20$ nodi.}
\label{fig:hb-cancellazione}
\end{figure}

\section{Analisi ammortizzata degli inserimenti}\label{sec:hb-ammortizzata}
Analizziamo ora il costo di $n$ inserimenti in un heap binomiale inizialmente vuoto.

Come costo effettivo dell'$i$-esimo inserimento consideriamo il numero di chiamate a \MakeBinomialHeap{} e a \BinomialLink{}, che costano entrambe $O(1)$:
\[
c_i = \#\MakeBinomialHeap + \#\BinomialLink = 2 + L_i,
\]
dove $2$ conta le due chiamate a \MakeBinomialHeap{}, una in \BinomialHeapInsert{} e una all'inizio di \BinomialHeapUnion{}, e $L_i$ è il numero di collegamenti eseguiti dall'$i$-esimo inserimento.

\paragraph{Metodo del potenziale.}
Definiamo come potenziale il numero di alberi binomiali dell'heap:
\[
\Phi(H) = \text{numero di alberi binomiali in } H.
\]
Indicando con $H_0$ l'heap iniziale vuoto, si ha $\Phi(H_0) = 0$ e $\Phi(H) \ge 0$ per ogni heap $H$. Per il lemma~\ref{lem:potenziale}, quindi, la somma dei costi ammortizzati è un limite superiore alla somma dei costi effettivi.

Calcoliamo il costo ammortizzato $\hat c = c + \Delta\Phi$ delle due parti dell'inserimento. La chiamata a \MakeBinomialHeap{} in \BinomialHeapInsert{} non cambia gli alberi di $H$, quindi:
\[
\hat c = c + \Delta\Phi = 1 + 0 = 1.
\]
L'unione, invece, esegue una chiamata a \MakeBinomialHeap{} e $L_i$ collegamenti. Il numero di alberi aumenta di uno per il nuovo $B_0$, e diminuisce di uno per ogni collegamento. Indicando con $T_i$ il numero di alberi dopo l'$i$-esimo inserimento si ha $T_i = T_{i-1} + 1 - L_i$, quindi:
\[
\begin{aligned}
\hat c &= c + \Delta\Phi = (1 + L_i) + (T_i - T_{i-1}) \\
&= (1 + L_i) + \big((T_{i-1} + 1 - L_i) - T_{i-1}\big) \\
&= 2.
\end{aligned}
\]
Il costo ammortizzato di un inserimento è quindi $\hat c_i = 1 + 2 = 3$, e $n$ inserimenti in un heap binomiale inizialmente vuoto costano:
\[
\sum_{i=1}^{n} c_i \le \sum_{i=1}^{n} \hat c_i = 3n = O(n).
\]
La figura~\ref{fig:hb-inserimenti} riporta i costi di otto inserimenti in un heap vuoto: il costo effettivo cambia da un inserimento all'altro, mentre quello ammortizzato vale sempre $3$.

\begin{figure}[htbp]
\centering
\begin{widget}{binomial-heap}[ops={insert 5, insert 3, insert 8, insert 1, insert 9, insert 2, insert 7, insert 4}, actions={insert}, potential, interval=1500]
\small
\begin{tabular}{cccccccc}
\toprule
$i$ & chiave & $n_i$ & $L_i$ & $c_i = 2 + L_i$ & $\Phi_i$ & $\Delta\Phi_i$ & $\hat c_i$ \\
\midrule
$1$ & 5 & $(0001)_2$ & $0$ & $2$ & $1$ & $+1$ & $3$ \\
$2$ & 3 & $(0010)_2$ & $1$ & $3$ & $1$ & $0$ & $3$ \\
$3$ & 8 & $(0011)_2$ & $0$ & $2$ & $2$ & $+1$ & $3$ \\
$4$ & 1 & $(0100)_2$ & $2$ & $4$ & $1$ & $-1$ & $3$ \\
$5$ & 9 & $(0101)_2$ & $0$ & $2$ & $2$ & $+1$ & $3$ \\
$6$ & 2 & $(0110)_2$ & $1$ & $3$ & $2$ & $0$ & $3$ \\
$7$ & 7 & $(0111)_2$ & $0$ & $2$ & $3$ & $+1$ & $3$ \\
$8$ & 4 & $(1000)_2$ & $3$ & $5$ & $1$ & $-2$ & $3$ \\
% totale c_i = 23
\midrule
 & & & & $\sum c_i = 23$ & & & $\sum \hat c_i = 24$ \\
\bottomrule
\end{tabular}
\end{widget}
\caption{Otto inserimenti in un heap binomiale inizialmente vuoto. Per ogni inserimento sono riportati la chiave, il numero di nodi $n_i$ in binario, i collegamenti $L_i$, il costo effettivo $c_i$, il potenziale $\Phi_i$ con la sua variazione, e il costo ammortizzato $\hat c_i$.}
\label{fig:hb-inserimenti}
\end{figure}

\paragraph{Metodo degli accantonamenti.}
Lo stesso risultato si ottiene con gli accantonamenti. Assegniamo alla chiamata a \MakeBinomialHeap{} in \BinomialHeapInsert{} un costo ammortizzato di $1$, che ne paga il costo, e all'unione un costo ammortizzato di $2$:
\begin{itemize}
    \item Un'unità paga la chiamata a \MakeBinomialHeap{} all'inizio dell'unione.
    \item L'altra unità viene assegnata al nuovo nodo, come credito per pagare il suo futuro collegamento.
\end{itemize}
Per induzione si verifica che la radice di ogni albero binomiale ha un'unità di credito. Il nuovo $B_0$ ha il credito appena assegnato. Quando due alberi vengono collegati, la radice che diventa figlio usa il suo credito per pagare \BinomialLink{}, mentre l'altra resta radice e conserva il proprio. In ogni istante, quindi, il credito totale è uguale al numero di alberi binomiali dell'heap, e non è mai negativo.

\section{Riepilogo dei costi}\label{sec:hb-costi}
La tabella~\ref{tab:hb-costi} riassume il costo nel caso peggiore delle procedure viste, e indica da dove nasce.

\begin{table}[htbp]
\centering
\small
\begin{tabular}{lcl}
\toprule
Procedura & Costo & Da dove nasce il costo \\
\midrule
\MakeBinomialHeap{} & $O(1)$ & allocazione dell'heap vuoto \\
\BinomialHeapMinimum{} & $O(\lg n)$ & scansione delle radici, al più $\lfloor \lg n \rfloor + 1$ \\
\BinomialLink{} & $O(1)$ & aggiornamento di quattro campi \\
\BinomialHeapMerge{} & $O(\lg n)$ & scansione delle due liste delle radici \\
\BinomialHeapUnion{} & $O(\lg n)$ & al più $2\lfloor \lg n \rfloor + 2$ radici nella lista fusa \\
\BinomialHeapInsert{} & $O(\lg n)$ & unione con un heap di un solo nodo \\
\BinomialHeapExtractMin{} & $O(\lg n)$ & ricerca del minimo, inversione dei figli e unione \\
\BinomialHeapDecreaseKey{} & $O(\lg n)$ & altezza degli alberi, al più $\lfloor \lg n \rfloor$ \\
\BinomialHeapDelete{} & $O(\lg n)$ & decremento della chiave ed estrazione del minimo \\
\bottomrule
\end{tabular}
\caption{Il costo nel caso peggiore delle procedure sugli heap binomiali, dove $n$ è il numero di nodi dell'heap. Su una sequenza di inserimenti in un heap vuoto, il numero di collegamenti per inserimento è costante in senso ammortizzato.}
\label{tab:hb-costi}
\end{table}

Rispetto agli heap binari, quindi, l'unione passa da $\Theta(n)$ a $O(\lg n)$, al prezzo di un minimo che costa $O(\lg n)$ invece di $\Theta(1)$. Tutte le operazioni restano logaritmiche, e su una sequenza di inserimenti in un heap vuoto il numero di collegamenti per inserimento è costante in senso ammortizzato.

\section*{Riferimenti}
Il materiale visto in questo capitolo comprende:
\begin{itemize}
    \item Materiale visto a lezione.
    \item Capitolo 19 (\emph{Binomial Heaps}) del libro T.\,H.~Cormen, C.\,E.~Leiserson, R.\,L.~Rivest, C.~Stein, \emph{Introduction to Algorithms}, seconda edizione.
    \item Introduzione del capitolo 19 (\emph{Fibonacci Heaps}) e problema 19-2 (\emph{Binomial trees and binomial heaps}) dello stesso libro, terza edizione.
\end{itemize}
